% PDF page 454
\source{gilbarg-trudinger:page-0454}

Here $F_{ij}(y)$ is the cofactor of $r_{ij}$, and $F$ is elliptic only for positive $r$. Accordingly,
(17.2) will be elliptic only for functions $u \in C^2(\Omega)$ that are uniformly convex at each
point of $\Omega$ and for such a solution to exist we must also have $f$ positive.

(ii) \emph{The Equation of Prescribed Gauss Curvature:} Let $u \in C^2(\Omega)$ and suppose the
graph of $u$ has Gauss curvature $K(x)$ at the point $(x,u(x))$, $x \in \Omega$. It follows (see
Section 14.6) that $u$ satisfies the equation

\begin{equation}
F[u] = \det D^2 u - K(x)(1 + |Du|^2)^{(n+2)/2} = 0,
\tag{17.3}
\end{equation}

which again will be elliptic only for uniformly convex $u \in C^2(\Omega)$. More generally,
examples (i) and (ii) may be combined into the family of equations of Monge-
Ampère type,

\begin{equation}
F[u] = \det D^2 u - f(x,u,Du) = 0,
\tag{17.4}
\end{equation}

where $f$ is a positive function on $\Omega \times \mathbb{R} \times \mathbb{R}^n$.

(iii) \emph{Pucci's Equations:} For $0 < \alpha \leq 1/n$, let $\mathcal{L}_\alpha$ denote the set of linear uniformly
elliptic operators of the form

\[
Lu = a^{ij}(x)D_{ij}u
\]

with bounded measurable coefficients $a^{ij}$ satisfying

\[
a^{ij}\xi_i \xi_j \geq \alpha |\xi|^2,\qquad \sum a^{ii} = 1
\]

for all $x \in \Omega, \xi \in \mathbb{R}^n$. Maximal and minimal operators $M_\alpha$ and $m_\alpha$ are then defined by

\begin{equation}
M_\alpha[u] = \sup_{L \in \mathcal{L}_\alpha} Lu,\qquad m_\alpha[u] = \inf_{L \in \mathcal{L}_\alpha} Lu.
\tag{17.5}
\end{equation}

The operators $M_\alpha$, $m_\alpha$ are fully nonlinear and are also related by $M_\alpha[-u] =
-m_\alpha[u]$. Furthermore, a simple calculation yields the formulae

\[
M_\alpha[u] = \alpha \Delta u + (1 - n\alpha)\mathcal{C}_n(D^2u),
\]
\[
m_\alpha[u] = \alpha \Delta u + (1 - n\alpha)\mathcal{C}_1(D^2u)
\]

where $\mathcal{C}_1(r), \mathcal{C}_n(r)$ are the minimum, maximum eigenvalues of the matrix $r$. We can
then consider the extremal equations

\begin{equation}
M_\alpha[u] = f;\qquad m_\alpha[u] = f
\tag{17.6}
\end{equation}

in the domain $\Omega$, for a given function $f$. Although the functions $\mathcal{C}_1, \mathcal{C}_n$ are not
differentiable, the concept of ellipticity readily extends to embrace this situation,
the equations (17.6) being in fact uniformly elliptic (see below).

(iv) \emph{The Bellman Equation:} When the family $\mathcal{L}_\alpha$ in the preceding example is
replaced by an arbitrary family $\mathcal{L}$ of linear elliptic operators we obtain the Bellman
equation for the optimal cost in a stochastic control problem; (see [KV 2]).

% PDF page 455
\source{gilbarg-trudinger:page-0455}

Specifically let us consider a family $\mathcal L$ indexed by a parameter $\nu$ belonging to a set
$V$. Suppose that each $L_\nu \in \mathcal L$ has the form

\[
(17.7)\qquad L_\nu u = a_\nu^{ij}(x)D_{ij}u + b_\nu^i(x)D_i u + c_\nu(x)u
\]

where $a_\nu^{ij}$, $b_\nu^i$, $c_\nu$ are real functions on $\Omega$ for each $i, j = 1, \ldots, n$, $\nu \in V$, and for each
$\nu \in V$, let $f_\nu$ be a real function on $\Omega$. The Bellman equation now takes the form

\[
(17.8)\qquad F[u] = \inf_{\nu \in V}(L_\nu u - f_\nu) = 0.
\]

There are various ways of extending the concept of ellipticity to nondifferentiable
$F$, for example by monotonicity or by approximation; (see Problem 17.1). For our
purposes it suffices to have an extension to functions $F$ which are Lipschitz con-
tinuous with respect to the $r$ variables. We then call the operator $F$ elliptic in a sub-
set $\mathcal U$ of $\Gamma$ if the matrix $[F_{ij}] = [\partial F/\partial r_{ij}]$ is positive wherever it exists in $\mathcal U$, and
uniformly elliptic in $\mathcal U$ if the ratio of maximum to minimum eigenvalues $\Lambda/\lambda$ is
bounded in $\mathcal U$. Note that $[F_{ij}]$ exists for almost all $r \in \mathbb R^{n \times n}$. With these definitions,
it follows that the Bellman equation (17.8) is elliptic in $\Omega$ if for each $x \in \Omega$, $\nu \in V$,

\[
(17.9)\qquad \lambda(x)|\xi|^2 \le a_\nu^{ij}(x)\xi_i\xi_j \le \Lambda(x)|\xi|^2
\]

for all $\xi \in \mathbb R^n$, where $\lambda$ and $\Lambda$ are positive functions in $\Omega$. Moreover, the Bellman
equation is uniformly elliptic in $\Omega$ if $\Lambda/\lambda \in L^\infty(\Omega)$.

\bigskip

\section*{17.1. Maximum and Comparison Principles}

The maximum and comparison principles derived in Chapter 10 for quasilinear
equations in general form are readily extended to fully nonlinear equations. We
shall establish the following form of the comparison principle.

\begin{theorem}[Theorem 17.1]\label{gt-ch17-thm-17-1}\dcref{17.1}\group{chapter-17}\source{gilbarg-trudinger:page-0455}
\emph{Let $u, v \in C^0(\overline{\Omega}) \cap C^2(\Omega)$ satisfy $F[u] \geq F[v]$ in $\Omega$, $u \leq v$ on $\partial \Omega$,
where:}

\begin{enumerate}
\item[(i)] \emph{the function $F$ is continuously differentiable with respect to the $z, p, r$ variables
in $\Gamma$;}
\item[(ii)] \emph{the operator $F$ is elliptic on all functions of the form $\theta u + (1 - \theta)v$,
$0 \leq \theta \leq 1$;}
\item[(iii)] \emph{the function $F$ is non-increasing in $z$ for each $(x, p, r) \in \Omega \times \mathbb R^n \times \mathbb R^{n \times n}$.}
\end{enumerate}

It then follows that $u \leq v$ in $\Omega$.
\end{theorem}

% PDF page 456
\source{gilbarg-trudinger:page-0456}

\textit{Proof.}\quad Writing

\[
w=u-v,
\]

\[
u_{\theta}=\theta u+(1-\theta)v,
\]

\[
a^{ij}(x)=\int_{0}^{1}F_{ij}(x,u_{\theta},Du_{\theta},D^{2}u_{\theta})\,d\theta,
\]

\[
b^{i}(x)=\int_{0}^{1}F_{p_{i}}(x,u_{\theta},Du_{\theta},D^{2}u_{\theta})\,d\theta,
\]

\[
c(x)=\int_{0}^{1}F_{z}(x,u_{\theta},Du_{\theta},D^{2}u_{\theta})\,d\theta,
\]

we obtain

\[
Lw=a^{ij}D_{ij}w+b^{i}D_{i}w+cw
=F[u]-F[v]\ge 0 \quad \text{in } \Omega.
\]

Furthermore conditions (i), (ii) and (iii) imply that $L$ satisfies the hypotheses of the
weak maximum principle (Theorem 3.1), and thus $w\le 0$ in $\Omega$. $\square$

\medskip

Weaker hypotheses are clearly possible in Theorem 17.1. Also, by virtue of the
strong maximum principle (Theorem 3.5), we have that either $u<v$ in $\Omega$ or $u$ and $v$
coincide. A uniqueness result for the Dirichlet problem follows immediately from
Theorem 17.1.

\medskip

\begin{corollary}[Corollary 17.2]\label{gt-ch17-cor-17-2}\dcref{17.2}\group{chapter-17}\source{gilbarg-trudinger:page-0456}
\textit{Let $u, v\in C^{0}(\bar{\Omega})\cap C^{2}(\Omega)$ satisfy $F[u]=F[v]$ in $\Omega$, $u=v$ on $\partial\Omega$
and suppose that conditions (i) to (iii) in Theorem 17.1 hold. Then $u\equiv v$ in $\Omega$.}
\end{corollary}

\medskip

Maximum principles, H\"older estimates for solutions and boundary gradient
estimates for fully nonlinear elliptic equations may often be inferred directly from
the corresponding results for quasilinear equations. If $u\in C^{2}(\Omega)$, we may write the
operator $F[u]$ in the form

\[
(17.10)\qquad F[u]=F(x,u,Du,D^{2}u)-F(x,u,Du,0)+F(x,u,Du,0)
\]
\[
\hspace{2.9cm}=a^{ij}(x,u,Du)D_{i,j}u+b(x,u,Du),
\]

where now

\[
a^{ij}(x,z,p)=\int_{0}^{1}F_{ij}(x,z,p,\theta D^{2}u(x))\,d\theta,
\]

\[
b(x,z,p)=F(x,z,p,0).
\]


% PDF page 457
\source{gilbarg-trudinger:page-0457}


In particular, defining
\[
\E(x, z, p, r) = F_{ij}(x, z, p, r)p_i p_j,
\qquad
\E^* = \E/|p|^2,
\]
\[
\D(x, z, p, r) = \det F_{ij}(x, z, p, r),
\qquad
\D^* = \D^{1/n},
\]
we thus obtain from Theorems 10.3 and 10.4 the following theorem.

\medskip

\begin{theorem}[Theorem 17.3]\label{gt-ch17-thm-17-3}\dcref{17.3}\group{chapter-17}\source{gilbarg-trudinger:page-0457}
\textit{Let $F$ be elliptic in $\Omega$ and suppose there exist non-negative constants $\mu_1$ and $\mu_2$ such that}

\begin{equation}\tag{17.11}
\frac{F(x, z, p, 0)\,\sign z}{\E^* \text{ (or } \D^*\text{)}} \leq \mu_1 |p| + \mu_2
\qquad \forall (x, z, p, r) \in \Gamma.
\end{equation}

\noindent\textit{Then, if $u \in C^0(\bar{\Omega}) \cap C^2(\Omega)$ satisfies $F[u] \geq 0$ $(=0)$ in $\Omega$, we have}

\begin{equation}\tag{17.12}
\sup_{\Omega} u(|u|) \leq \sup_{\partial\Omega} u^+(|u|) + C\mu_2
\end{equation}

\noindent\textit{where $C = C(\mu_1, \operatorname{diam}\Omega)$.}
\end{theorem}

\medskip

For equations of Monge-Amp\`ere type a maximum principle analogous to Theorem 10.5 can be concluded directly from Lemma 9.4.

\medskip

\begin{theorem}[Theorem 17.4]\label{gt-ch17-thm-17-4}\dcref{17.4}\group{chapter-17}\source{gilbarg-trudinger:page-0457}
\textit{Let $F$ be given by (17.4) and suppose there exist non-negative functions $g \in L^1_{\mathrm{loc}}(\mathbb{R}^n)$, $h \in L^1(\Omega)$ and a constant $N$ such that}

\begin{equation}\tag{17.13}
\lvert f(x, z, p)\rvert \leq \frac{h(x)}{g(p)}
\qquad \forall x \in \Omega,\ |z| \geq N,\ p \in \mathbb{R}^n;
\end{equation}

\begin{equation}\tag{17.14}
\int_{\Omega} h\,dx < \int_{\mathbb{R}^n} g(p)\,dp \equiv g_{\infty}.
\end{equation}

\noindent\textit{Then, if $u \in C^0(\bar{\Omega}) \cap C^2(\Omega)$ satisfies $F[u] \geq 0$ $(=0)$ in $\Omega$, we have}

\begin{equation}\tag{17.15}
\sup_{\Omega} u(|u|) \leq \sup_{\partial\Omega} u^+(|u|) + C\,\operatorname{diam}\Omega + N,
\end{equation}

\noindent\textit{where $C$ depends only on $g$ and $h$. In particular, if $F$ is given by (17.2), we obtain}

\begin{equation}\tag{17.16}
\sup_{\Omega} u(|u|) \leq \sup_{\partial\Omega} u^+(|u|) + \frac{\operatorname{diam}\Omega}{(\omega_n)^{1/n}}\left(\int_{\Omega} |f|\right)^{1/n};
\end{equation}

\noindent\textit{while if $F$ is given by (17.3), the estimate (17.15) holds with $C = C(n, K_0)$, provided}

\begin{equation}\tag{17.17}
K_0 = \int_{\Omega} \lvert K(x)\rvert < \omega_n.
\end{equation}
\end{theorem}

% PDF page 458
\source{gilbarg-trudinger:page-0458}

\begin{quote}
To conclude this section we note that the preceding results and their proofs also extend to nondifferentiable functions $F$. In particular, we obtain as a generalization of the comparison principle (Theorem 17.1):
\end{quote}

\begin{theorem}[Theorem 17.5]\label{gt-ch17-thm-17-5}\dcref{17.5}\group{chapter-17}\source{gilbarg-trudinger:page-0458}
\textit{Let $u, v \in C^0(\bar{\Omega}) \cap C^2(\Omega)$ satisfy $F[u] \geq F[v]$ in $\Omega$, $u \leq v$ on $\partial\Omega$, where}

\begin{enumerate}
\item[(i)] $F$ is locally uniformly Lipschitz with respect to the $z, p, r$ variables in $\Gamma$;
\item[(ii)] $F$ is elliptic in $\Omega$;
\item[(iii)] $F$ is non-increasing in $z$ for each $(x, p, r) \in \Omega \times \mathbb{R}^n \times \mathbb{R}^{n \times n}$;
\item[(iv)] $\lvert F_p\rvert/\lambda$ is locally bounded in $\Gamma$.
\end{enumerate}

\noindent\textit{It then follows that $u \leq v$ in $\Omega$.}
\end{theorem}

\medskip

As an application of Theorem 17.5 we see that the comparison principle will hold for the Bellman operators (17.8) provided (17.9) holds, $c_\nu \leq 0$ for all $\nu \in V$ and $\lvert b_\nu\rvert/\lambda$ is locally bounded in $\Omega$, uniformly in $\nu$. By virtue of the maximum principle for strong solutions (Theorem 9.1), all the results of this section carry over to functions $u, v \in W^{2,n}_{\mathrm{loc}}(\Omega)$.

\section*{17.2. The Method of Continuity}

The topological methods of Chapter 11 are inadequate for the treatment of fully nonlinear elliptic equations or nonlinear boundary value problems. For these problems, we shall use a nonlinear version of the method of continuity (Theorem 5.2). In principle, the continuity method involves the embedding of the given problem in a family of problems indexed by a bounded closed interval, say $[0, 1]$. The subset $S$ of $[0, 1]$ for which the corresponding problems are solvable is shown to be nonempty, closed and open, and hence it coincides with the whole interval. As in the quasilinear case in Chapter 11, the linear theory is again vital but in the present situation it is applied to the Fréchet derivative of the operator $F$ in order to demonstrate the openness of the solvability set $S$.

We commence with an abstract functional analytic formulation. Let $\mathcal{B}_1$ and $\mathcal{B}_2$ be Banach spaces and $F$ a mapping from an open set $U \in \mathcal{B}_1$ into $\mathcal{B}_2$. The mapping $F$ is called \textit{Fr\'echet differentiable} at an element $u \in \mathcal{B}_1$ if there exists a bounded linear mapping $L : \mathcal{B}_1 \to \mathcal{B}_2$ such that

\begin{equation}\tag{17.18}
\frac{\lVert F[u + h] - F[u] - Lh\rVert_{\mathcal{B}_2}}{\lVert h\rVert_{\mathcal{B}_1}} \to 0
\end{equation}

\noindent as $h \to 0$ in $\mathcal{B}_1$. The linear mapping $L$ is called the \textit{Fr\'echet derivative} (or \textit{differential}) of $F$ at $u$ and will be denoted by $F_u$. When $\mathcal{B}_1, \mathcal{B}_2$ are Euclidean spaces, $\mathbb{R}^n, \mathbb{R}^m$, the Fr\'echet derivative coincides with the usual notion of differential, and, moreover, the basic theory for the infinite dimensional case can be modelled on that for the finite dimensional case as usually treated in advanced calculus; (see for example [DI]). In particular, it is evident from (17.18) that the Fr\'echet differentiability of $F$ at $u$

% PDF page 459
\source{gilbarg-trudinger:page-0459}

\pagestyle{empty}



implies that $F$ is continuous at $u$ and that the Fr\'echet derivative $F_u$ is determined uniquely by (17.18). We call $F$ continuously differentiable at $u$ if $F$ is Fr\'echet differentiable in a neighbourhood of $u$ and the resulting mapping

\[
v \mapsto F_v \in \E(\B_1,\B_2)
\]

is continuous at $u$. Here $\E(\B_1,\B_2)$ denotes the Banach space of bounded linear mappings from $\B_1$ into $\B_2$ with norm given by

\[
\norm{L} = \sup_{v \in \B_1 \atop v \neq 0} \frac{\norm{Lv}_{\B_2}}{\norm{v}_{\B_1}}.
\]

The chain rule holds for Fr\'echet differentiation, that is, if $F : \B_1 \to \B_2$, $G : \B_2 \to \B_3$ are Fr\'echet differentiable at $u \in \B_1$, $F[u] \in \B_2$, respectively, then the composite mapping $G \circ F$ is differentiable at $u \in \B_1$ and

\begin{equation}
(G \circ F)_u = G_{F[u]} \circ F_u
\tag{17.19}
\end{equation}

The theorem of the mean also holds in the sense that if $u, v \in \B_1$, $F : \B_1 \to \B_2$ is differentiable on the closed line segment $\gamma$ joining $u$ and $v$ in $\B_1$, then

\begin{equation}
\norm{F[u] - F[v]}_{\B_2} \leq K \norm{u - v}_{\B_1},
\tag{17.20}
\end{equation}

where

\[
K = \sup_{w \in \gamma} \norm{F_w}.
\]

With the aid of these basic properties we may deduce an implicit function theorem for Fr\'echet differentiable mappings. Suppose that $\B_1$, $\B_2$ and $X$ are Banach spaces and that $G : \B_1 \times X \to \B_2$ is Fr\'echet differentiable at a point $(u,\sigma)$, $u \in \B_1$, $\sigma \in X$. The partial Fr\'echet derivatives, $G^{1}_{(u,\sigma)}, G^{2}_{(u,\sigma)}$ at $(u,\sigma)$ are the bounded linear mappings from $\B_1$, $X$, respectively, into $\B_2$ defined by

\[
G_{(u,\sigma)}(h,k) = G^{1}_{(u,\sigma)}(h) + G^{2}_{(u,\sigma)}(k)
\]

for $h \in \B_1$, $k \in X$. We state the implicit function theorem in the following form.

\begin{theorem}[Theorem 17.6]\label{gt-ch17-thm-17-6}\dcref{17.6}\group{chapter-17}\source{gilbarg-trudinger:page-0460}
\textit{Let $\B_1$, $\B_2$ and $X$ be Banach spaces and $G$ a mapping from an open subset of $\B_1 \times X$ into $\B_2$. Let $(u_0,\sigma_0)$ be a point in $\B_1 \times X$ satisfying:}

\begin{enumerate}
\item[(i)] $G[u_0,\sigma_0] = 0$;
\item[(ii)] $G$ is continuously differentiable at $(u_0,\sigma_0)$;
\item[(iii)] the partial Fr\'echet derivative $L = G^{1}_{(u_0,\sigma_0)}$ is invertible.
\end{enumerate}

\textit{Then there exists a neighbourhood $\mathcal{N}$ of $\sigma_0$ in $X$ such that the equation $G[u,\sigma] = 0$, is solvable for each $\sigma \in \mathcal{N}$, with solution $u = u_\sigma \in \B_1$.}
\end{theorem}


% PDF page 460
\source{gilbarg-trudinger:page-0460}


Theorem 17.6 may be proved by reduction to the contraction mapping principle
(Theorem 5.1). Indeed the equation, $G[u,\sigma]=0$, is equivalent to the equation

\[
u=Tu\equiv u-L^{-1}G[u,\sigma],
\]

and the operator $T$ will, by virtue of (17.19) and (17.20), be a contraction mapping
in a closed ball $\bar{B}=\bar{B}_{\delta}(u_{0})$ in $\B_{1}$ for $\delta$ sufficiently small and $\sigma$ sufficiently close to $\sigma_{0}$
in $X$; (see [DI] for details). One can further show that the mapping $F : X \to \B_{1}$
defined by $\sigma\to u_{\sigma}$ for $\sigma\in\N$, $u_{\sigma}\in\bar{B}$, $G[u_{\sigma},\sigma]=0$ is differentiable at $\sigma_{0}$ with Fréchet
derivative

\[
F_{\sigma_{0}}=-L^{-1}G^{2}_{(u_{0},\sigma_{0})}.
\]

In order to apply Theorem 17.6 we suppose that $\B_{1}$ and $\B_{2}$ are Banach spaces
with $F$ a mapping from an open subset $\U \subset \B_{1}$ into $\B_{2}$. Let $\psi$ be a fixed element in
$\U$ and define for $u\in\U$, $t\in\mathbb{R}$ the mapping $G : \U \times \mathbb{R} \to \B_{2}$ by

\[
G[u,t]=F[u]-tF[\psi].
\]

Let $S$ and $E$ be the subsets of $[0,1]$ and $\B_{1}$ defined by

\begin{equation}
\tag{17.21}
\begin{aligned}
S &= \{t\in[0,1]\mid G[u,t]=0 \quad \text{for some }u\in\U\},\\
E &= \{u\in\U\mid G[u,t]=0 \quad \text{for some }t\in[0,1]\}.
\end{aligned}
\end{equation}

Clearly $1\in S$, $\psi\in E$ so that $S$ and $E$ are not empty. Let us next suppose that the
mapping $F$ is continuously differentiable on $E$ with invertible Fréchet derivative $F_{u}$.
It follows then from the implicit function theorem (Theorem 17.6) that the set $S$ is
open in $[0,1]$. Consequently we obtain the following version of the method of
continuity.

\begin{theorem}[Theorem 17.7]\label{gt-ch17-thm-17-7}\dcref{17.7}\group{chapter-17}\source{gilbarg-trudinger:page-0460}
\textit{The equation $F[u]=0$ is solvable for $u\in\U$ provided the set $S$ is
closed in $[0,1]$.}
\end{theorem}

Let us now examine the application of Theorem 17.7 to the Dirichlet problem for
fully nonlinear elliptic equations. We assume that the function $F$ in equation (17.1)
belongs to $C^{2,\alpha}(\bar{\Gamma})$ and for the Banach spaces $\B_{1}$ and $\B_{2}$ take $\B_{1}=C^{2,\alpha}(\bar{\Omega})$,
$\B_{2}=C^{\alpha}(\bar{\Omega})$ for some $\alpha\in(0,1)$. Clearly the operator $F$ defined by (17.1) maps $\B_{1}$
into $\B_{2}$, and furthermore by calculation we see that $F$ has continuous Fréchet
derivative $F_{u}$ given by

\begin{equation}
\tag{17.22}
F_{u}h=Lh=F_{ij}(x)D_{ij}h+b^{i}(x)D_{i}h+c(x)h,
\end{equation}

% PDF page 461
\source{gilbarg-trudinger:page-0461}

17.2. The Method of Continuity

where

\[
F_{ij}(x)=F_{ij}(x,u(x),Du(x),D^2u(x)),
\]
\[
b^i(x)=F_{p_i}(x,u(x),Du(x),D^2u(x)),
\]
\[
c(x)=F_z(x,u(x),Du(x),D^2u(x));
\]

(see Problem 17.2). We cannot expect the mapping $F_u$ to be invertible on all of
$C^{2,\alpha}(\overline{\Omega})$ and we accordingly restrict $F$ to the subspace $\B_1
=\{u\in C^{2,\alpha}(\overline{\Omega})\mid u=0$ on
$\partial\Omega\}$, (already used in Section 6.3). The linear operator $L$ will then be invertible for
any $u\in C^{2,\alpha}(\overline{\Omega})$ for which $L$ is strictly elliptic and $c\le 0$ in $\Omega$, provided $\partial\Omega\in C^{2,\alpha}$
(Theorem 6.14). By means of Theorem 17.7 the solvability of the Dirichlet problem
is now reduced to the establishment of apriori estimates in the space $C^{2,\alpha}(\overline{\Omega})$.

\medskip

\noindent\textbf{Theorem 17.8.} \textit{Let $\Omega$ be a bounded domain in $\R^n$ with boundary $\partial\Omega\in C^{2,\alpha}$, $0<\alpha<1$, $\U$ an open subset of the space $C^{2,\alpha}(\overline{\Omega})$ and $\phi$ a function in $\U$. Set $E=\{u\in\U\mid F[u]=\sigma F[\phi]$ for some $\sigma\in[0,1],\ u=\phi$ on $\partial\Omega\}$ and suppose that $F\in C^{2,\alpha}(\overline{\Gamma})$ together with}

\begin{enumerate}
\item[(i)] $F$ is strictly elliptic in $\Omega$ for each $u\in E$;
\item[(ii)] $F_z(x,u,Du,D^2u)\le 0$ for each $u\in E$;
\item[(iii)] $E$ is bounded in $C^{2,\alpha}(\overline{\Omega})$;
\item[(iv)] $\overline{E}\subset \U$.
\end{enumerate}

\noindent \textit{Then the Dirichlet problem, $F[u]=0$ in $\Omega$, $u=\phi$ on $\partial\Omega$ is solvable in $\U$.}

\medskip

\noindent\textbf{Proof.} We can reduce to the case of zero boundary values by replacing $u$ with
$u-\phi$. The mapping $G:\B_1\times \R\to \B_2$ is then defined by taking $\psi\equiv 0$ so that

\[
G[u,\sigma]=F[u+\phi]-\sigma F[\phi].
\]

It then follows from Theorem 17.7 and the remarks preceding the statement of
Theorem 17.8, that the given Dirichlet problem is solvable provided the set $S$ is
closed. However the closure of $S$ (and also $E$) follows readily from the boundedness
of $E$ (and hypothesis (iv)) by virtue of the Arzela-Ascoli theorem. $\Box$

\medskip

When specialized to the case of quasilinear equations, Theorem 17.7 is somewhat weaker than Theorem 11.8. For the quasilinear case, estimates in $C^{1,\beta}(\overline{\Omega})$ for
some $\beta\in(0,1)$ will imply, by the Schauder theory, estimates in $C^{2,\alpha}(\overline{\Omega})$. The requirement $F\in C^2(\overline{\Gamma})$ can be weakened (Problem 17.3), but in order to assure the desired Fr\'echet differentiability of $F$ we still need smoother coefficient hypotheses than those of Theorem 11.8.

Theorem 17.8 reduces the solvability of the Dirichlet problem for fully nonlinear elliptic equations to a series of derivative estimates extending Steps I to IV in the existence procedure for the quasilinear case as described in Chapter 11. The


% PDF page 462
\source{gilbarg-trudinger:page-0462}

\pagestyle{empty}


fully nonlinear problem is in general more involved since the estimation of second
derivatives is required. In the ensuing sections we shall establish second derivative
estimates for various types of equations including the examples mentioned at the
beginning of this chapter. In certain cases, these estimates will be insufficient for the
direct application of Theorem 17.7 and some modifications will be necessary. In
particular, the nonsmoothness of the function $F$ in the Bellman equation (17.8) is
overcome through approximation. For uniformly elliptic equations we shall take
$\mathcal{U} = C^{2,\alpha}(\bar{\Omega})$ (so that hypothesis (iv) becomes redundant) while for equations of
Monge-Ampere type, $\mathcal{U}$ will be the subset of uniformly convex functions and hy-
pothesis (iv) is ensured by the positivity of the function $f$. While the family of
operators,

\[
F[u; \sigma] = F[u] - \sigma F[\phi], \qquad \sigma \in [0,1]
\]

suffices for the applications in this chapter, it evidently can be replaced by any
family

\[
F[u; \sigma] = F(x, u, D u, D^2 u; \sigma)
\]

with $F \in C^2(\bar{\Gamma} \times [0,1])$, $F[u;0] = F[u]$, for which the Dirichlet problem,
$F[u;1] = 0$ in $\Omega$, $u = \phi$ on $\partial\Omega$, is solvable.

\section*{17.3. Equations in Two Variables}

For fully nonlinear equations in two variables, the H\"older gradient estimates in
Sections 12.2 and 13.2 enable the a priori estimation required for hypothesis (iii)
in Theorem 17.8 to be reduced to estimation in $C^2(\bar{\Omega})$. To show this, we assume that
$u \in C^3(\Omega)$ is a solution of (17.1) in $\Omega$ and differentiate with respect to the variable $x_k$,
thereby obtaining the equation

\[
(17.23) \qquad F_{ij} D_{ijk}u + F_{p_i} D_{ik}u + F_z D_k u + F_{x_k} = 0,
\]

where the arguments of the partial derivatives $F_{ij}, F_{p_i}, F_z, F_{x_k}$ are $x, u, D u, D^2u$.
Setting $w = D_k u$, $f = F_{x_k}(x, u, D u, D^2u)$, we may write (17.23), in the notation of
the preceding section, as

\[
Lw = a^{ij}D_{ij}w + b^iD_i w + c w = -f.
\]

Consequently, if $F$ is elliptic with respect to $u$, the first derivatives of $u$ will be solu-
tions of linear elliptic equations in $\Omega$. Accordingly, taking $n = 2$, and letting $\lambda, \Lambda, \mu$
satisfy

\[
(17.24) \qquad 0 < \lambda |\xi|^2 \leq F_{ij}(x, u, D u, D^2u)\xi_i\xi_j \leq \Lambda |\xi|^2,
\]
\[
|F_p, F_z, F_x(x, u, D u, D^2u)| \leq \lambda\mu,
\]


% PDF page 463
\source{gilbarg-trudinger:page-0463}


\noindent for all nonzero $\xi \in \mathbb{R}^{2}$ we have from Theorem 12.4 or 13.3:

\medskip

\noindent \textbf{Theorem 17.9.} \textit{Let $u \in C^{3}(\Omega)$ satisfy $F[u] = 0$ in $\Omega \subset \mathbb{R}^{2}$, where $F \in C^{1}(\Gamma)$, $F$ is elliptic with respect to $u$ and (17.24) is satisfied. Then for any $\Omega' \Subset \Omega$, we have the estimate}

\[
(17.25)\qquad [D^{2}u]_{\alpha;\Omega'} \le \frac{C}{d^{\alpha}}\{\,|D^{2}u|_{0;\Omega} + \mu d(1 + |Du|_{1;\Omega})\,\}
\]

\noindent \textit{where $C$ and $\alpha$ depend only on $\Lambda/\lambda$, and $d = \operatorname{dist}(\Omega', \partial\Omega)$.}

\medskip

\noindent To derive the corresponding global estimate we return initially to the general case $n \ge 2$ and assume that $\partial\Omega \in C^{3}$, $u \in C^{3}(\Omega) \cap C^{2}(\bar{\Omega})$ and $u = \phi$ on $\partial\Omega$ where $\phi \in C^{3}(\bar{\Omega})$. Let us fix a point $x_{0} \in \partial\Omega$ and flatten the part of $\partial\Omega$ in a ball $B = B(x_{0})$, centered at $x_{0}$, by a one-to-one mapping $\psi$ from $B$ onto an open set $D \subset \mathbb{R}^{n}$ such that

\[
\psi(B \cap \Omega) \subset \mathbb{R}^{n}_{+} = \{x \in \mathbb{R}^{n} \mid x_{n} > 0\},
\]
\[
\psi(B \cap \partial\Omega) \subset \partial\mathbb{R}^{n}_{+} = \{x \in \mathbb{R}^{n} \mid x_{n} = 0\},
\]
\[
\psi \in C^{3}(B),\qquad \psi^{-1} \in C^{3}(D).
\]

\noindent Writing

\[
 y = \psi(x),\qquad w = D_{y_{k}}u = \frac{\partial x_{l}}{\partial y_{k}} D_{l}u
\]

\noindent we then have

\[
 w = D_{y_{k}}\phi \quad \text{on } B \cap \partial\Omega,\qquad k = 1, \ldots, n-1
\]

\noindent and, using (17.23), we obtain in $B \cap \Omega$ the equation

\[
(17.26)\qquad F_{ij}D_{ij}w + F_{p_{i}}D_{i}w + F_{z}w
\]
\[
\qquad = 2F_{ij}D_{i}\!\left(\frac{\partial x_{l}}{\partial y_{k}}\right)D_{j,l}u + F_{ij}D_{ij}\!\left(\frac{\partial x_{l}}{\partial y_{k}}\right)D_{l}u + F_{p_{i}}D_{i}\!\left(\frac{\partial x_{l}}{\partial y_{k}}\right)D_{l}u - F_{x_{l}}\!\left(\frac{\partial x_{l}}{\partial y_{k}}\right).
\]

\noindent Consequently, if $n = 2$ and the hypotheses of Theorem 17.9 are satisfied, we conclude, from the arguments of Sections 12.1 and 12.2 or 13.1 and 13.2, a Hölder estimate for $Dw$ in a neighbourhood of $x_{0}$, for the case $k = 1. $ Furthermore, for any $D' \Subset D$, $y_{0} \in D^{+} \cap D'$, $R \le d/3$, where $D^{+} = \psi(B \cap \Omega)$, $d = \operatorname{dist}(D', \partial D)$, we have

\[
\int_{B_{R} \cap D^{+}} |D_{y}^{2}w|^{2}\,dy \le CR^{n-2+2\alpha}\{(1 + \mu)(1 + |Du|_{1;\Omega} + |D^{3}\phi|_{0;\Omega})^{2},
\]


% PDF page 464
\source{gilbarg-trudinger:page-0464}




where $\alpha=\alpha(\Lambda/\lambda,\psi)>0$ and $C=C(\Lambda/\lambda,\psi,\diam\Omega)$. Using equation (17.26) in the case $k=2$, we then obtain

\begin{equation}
\int_{B_R\cap D^+} \left|D_y^3u\right|^2\,dy \le CR^{n-2+2\alpha}\{(1+\mu)(1+|Du|_{1;\Omega})+|D^3\phi|_{0;\Omega}^2\}
\tag{17.27}
\end{equation}

and the desired global H\"older estimate follows by Morrey's estimate, Theorem 7.19.

\textbf{Theorem 17.10.} \textit{Let $u\in C^3(\Omega)\cap C^2(\overline{\Omega})$ satisfy $F[u]=0$ in $\Omega\subset \mathbb{R}^2$, $u=\phi$ on $\partial\Omega$ where $F\in C^1(\Gamma)$, $F$ is elliptic with respect to $u$, $\partial\Omega\in C^3$, $\phi\in C^3(\overline{\Omega})$ and (17.24) is satisfied. Then we have the estimate}

\begin{equation}
[D^2u]_{\alpha;\Omega}\le C\{(1+\mu)(1+|u|_{2;\Omega})+|\phi|_{3;\Omega}\}
\tag{17.28}
\end{equation}

\textit{where $\alpha$ and $C$ depend only on $\Lambda/\lambda$ and $\Omega$.}

As a consequence of Theorem 17.10 (and the regularity result, Lemma 17.16) the space $C^{2,\alpha}(\overline{\Omega})$ in hypothesis (iii) of Theorem 17.8 can be replaced by $C^2(\overline{\Omega})$, provided $\partial\Omega\in C^3$ and $\phi\in C^3(\overline{\Omega})$. For certain equations the necessary estimation of second derivatives can be achieved through interpolation, as in the estimation of first derivatives in Chapter 12. Indeed, let us assume the following structure conditions:

\begin{equation}
\begin{aligned}
0<\lambda|\xi|^2 &\le F_{ij}(x,z,p,r)\xi_i\xi_j \le \Lambda|\xi|^2 \\
\left|F_p,F_z(x,z,p,r)\right| &\le \mu\lambda, \\
\left|F_x(x,z,p,r)\right| &\le \mu\lambda(1+|p|+|r|).
\end{aligned}
\tag{17.29}
\end{equation}

for all nonzero $\xi\in\mathbb{R}^2$, $(x,z,p,r)\in\Gamma$, where $\lambda$ is a nonincreasing function of $|z|$, and $\Lambda$ and $\mu$ are nondecreasing functions of $|z|$. The estimates (17.25) and (17.28) will then be valid with $\lambda=\lambda(M)$, $\Lambda=\Lambda(M)$, $\mu=\mu(M)$ where $M=|u|_{0;\Omega}$. Consequently, with the aid of the interpolation inequalities, Lemmas 6.32 and 6.35, we can estimate the norms $|u|_{2,\alpha;\Omega}^*$, $|u|_{2,\alpha;\Omega}$ in terms of $M$.

\textbf{Theorem 17.11.} \textit{Let $u\in C^3(\Omega)$ satisfy $F[u]=0$ in $\Omega\subset\mathbb{R}^2$ and assume the structure conditions (17.29). Then we have the interior estimate}

\begin{equation}
|u|_{2,\alpha;\Omega}^*\le C
\tag{17.30}
\end{equation}

\textit{where $\alpha>0$ depends only on $\Lambda/\lambda$ and $C$ depends on $\Lambda/\lambda$, $\mu$ and $|u|_{0;\Omega}$. In addition, if $u\in C^3(\overline{\Omega})$, $\partial\Omega\in C^3$ and $u=\phi$ on $\partial\Omega$ for $\phi\in C^3(\overline{\Omega})$, we have the global estimate}

\begin{equation}
|u|_{2,\alpha;\Omega}\le C
\tag{17.31}
\end{equation}

\textit{where $\alpha>0$ depends on $\Lambda/\lambda$ and $\Omega$, and $C$ depends on $\Lambda/\lambda$, $\mu$, $\Omega$ and $|u|_{0;\Omega}$.}


% PDF page 465
\source{gilbarg-trudinger:page-0465}

We remark here that the hypotheses of Theorems 17.8, 17.9 and 17.10 can be weakened to permit
$F \in C^{0,1}(\Gamma)$, $u \in W^{3,2}(\Omega)$ and $\phi \in W^{3,2}(\Omega) \cap C^{2,\beta}(\overline{\Omega})$, (with the structure conditions (17.24), (17.28) holding almost everywhere in $\Omega$, $\Gamma$, respectively). The differentiation of equation (17.1) is accomplished with the aid of a generalization of the chain rule (Theorem 7.8). The structure conditions (17.29) may also be relaxed to correspond with the natural conditions for quasilinear equations; (see Notes).

By combining Theorems 17.3, 17.8 and 17.11, we obtain an existence theorem for the Dirichlet problem.

\textbf{Theorem 17.12.} \textit{Let $\Omega$ be a bounded domain in $\mathbb{R}^2$ with boundary $\partial\Omega \in C^3$ and let $\phi \in C^3(\overline{\Omega})$. Suppose that the operator $F$ satisfies $F_z \leq 0$ in $\Gamma$ together with the structure condition (17.29). Then the classical Dirichlet problem, $F[u] = 0$, $u = \phi$ on $\partial\Omega$, is uniquely solvable, with solution $u \in C^{2,\alpha}(\overline{\Omega})$ for all $\alpha < 1$.}

Note that for the direct application of Theorem 17.8 we require smoother $F$; the full strength of Theorem 17.12 then follows by approximation of $F$ and the estimate (17.31). Similar approximation also yields an existence theorem for equations of Bellman-Pucci type which we take up in Section 17.5.

\section*{17.4. Hölder Estimates for Second Derivatives}

We derive in this section interior Hölder estimates for second derivatives of solutions of fully nonlinear elliptic equations under the key assumption that the function $F$ is a concave (or convex) function of the $r$ variables. This restriction, which was not necessary in the two variable case of the preceding section, still enables us to cover the equations of Monge-Ampère and Bellman-Pucci type. To illustrate the main features of the method, we first consider equations of the special form

\begin{equation}
(17.32) \qquad F(D^2u) = g,
\end{equation}

where $F \in C^2(\mathbb{R}^{n \times n})$, $g \in C^2(\Omega)$ and $u \in C^4(\Omega)$. We suppose:

\begin{enumerate}
\item[(i)] $F$ is uniformly elliptic with respect to $u$, so that there exist positive constants $\lambda, \Lambda$ such that
\end{enumerate}

\begin{equation}
(17.33) \qquad \lambda |\xi|^2 \leq F_{ij}(D^2u)\xi_i\xi_j \leq \Lambda |\xi|^2
\end{equation}

for all $\xi \in \mathbb{R}^n$;
\begin{enumerate}
\item[(ii)] $F$ is concave with respect to $u$ in the sense that $F$ is a concave function on the range of $D^2u$.
\end{enumerate}

% PDF page 466
\source{gilbarg-trudinger:page-0466}

Let $\gamma$ be an arbitrary unit vector in $\mathbb{R}^n$ and differentiate equation (17.32) twice in the direction $\gamma$. We thus obtain

\begin{equation}\tag{17.34}
F_{ij}D_{ij\gamma}u = D_{\gamma}g,
\end{equation}
\[
F_{ij}D_{ij\gamma\gamma}u + F_{ij,kl}D_{ij\gamma}uD_{kl\gamma}u = D_{\gamma\gamma}g,
\]
where
\[
[F_{ij,kl}] = \left[\frac{\partial^2F}{\partial r_{ij}\,\partial r_{kl}}\right]
\]
is nonpositive by virtue of the concavity of $F$. Consequently, the function $w = D_{\gamma\gamma}u$ satisfies the differential inequality

\begin{equation}\tag{17.35}
F_{ij}D_{ij}w \ge D_{\gamma\gamma}g
\end{equation}
in $\Omega$. We now invoke the weak Harnack inequality from Section 9.7. Let $B_R$, $B_{2R}$ be concentric balls in $\Omega$ of radii $R$, $2R$, respectively, and set for $s = 1,2$,

\[
M_s = \sup_{B_{sR}} w, \qquad m_s = \inf_{B_{sR}} w.
\]

Applying Theorem 9.22 to the function $M_2 - w$, we thus obtain

\begin{equation}\tag{17.36}
\left\{R^{-n}\int_{B_R}(M_2 - w)^p\right\}^{1/p}
\le C\{M_2 - M_1 + R\|(D_{\gamma\gamma}g)^-\|_{n;B_{2R}}\}
\le C\{M_2 - M_1 + R^2|D^2g|_{0;\Omega}\}
\end{equation}

To conclude a Hölder estimate for $w$ from (17.36), we need a corresponding inequality for $-w$, which we obtain by considering (17.32) as a functional relation-ship between the second derivatives of $u$. To begin with, using the concavity of $F$ again, we have for any $x$, $y\in\Omega$,

\begin{equation}\tag{17.37}
F_{ij}(D^2u(y))(D_{ij}u(x) - D_{ij}u(y)) \ge F(D^2u(x)) - F(D^2u(y))
= g(x) - g(y).
\end{equation}

We now get a relationship between pure second derivatives of $u$ by means of the following matrix result (from [MW]).

\textbf{Lemma 17.13.} \textit{Let $S[\lambda, \Lambda]$ denote the set of positive matrices in $\mathbb{R}^{n\times n}$ with eigen-values lying in the interval $[\lambda, \Lambda]$, where $0 < \lambda < \Lambda$. Then there exists a finite set of unit vectors, $\gamma_1,\ldots,\gamma_N \in \mathbb{R}^n$ and positive numbers $\lambda^* < \Lambda^*$, depending only on $n$, $\lambda$}

% PDF page 467
\source{gilbarg-trudinger:page-0467}

\textit{and $\Lambda$ such that any matrix $\mathcal A=[a^{ij}]\in S[\lambda,\Lambda]$ can be written in the form}

\[
(17.38)\qquad \mathcal A=\sum_{k=1}^{N}\beta_k\gamma_k\otimes\gamma_k,\qquad \text{i.e.}\quad
a^{ij}=\sum_{k=1}^{N}\beta_k\gamma_{ki}\gamma_{kj},
\]

\textit{where $\lambda^{*}\leq \beta_k\leq \Lambda^{*}$, $k=1,\ldots,N$. Furthermore the directions
$\gamma_1,\ldots,\gamma_N$ can be chosen to include the coordinate directions $e_i$, $i=1,\ldots,n$, together with the directions
$(1/\sqrt{2})(e_i\pm e_j)$, $i<j$, $i,j=1,\ldots,n$.}

The proof of Lemma 17.13 is deferred until the end of this section. Applying Lemma 17.13 to the matrix $F_{ij}$, we obtain from (17.37) the inequality

\[
(17.39)\qquad \sum_{k=1}^{N}\beta_k\bigl(w_k(y)-w_k(x)\bigr)\leq g(y)-g(x)
\]

where $w_k=D_{\gamma_k\gamma_k}u$ and $\beta_k=\beta_k(y)$ satisfy

\[
0<\lambda^{*}\leq \beta_k\leq \Lambda^{*},
\]

the vectors $\gamma_1,\ldots,\gamma_N$ and numbers $\lambda^{*},\Lambda^{*}$ depending only on $n$, $\lambda$, $\Lambda$. Setting

\[
M_{sk}=\sup_{B_{sR}} w_k,\qquad m_{sk}=\inf_{B_{sR}} w_k,\qquad s=1,2;\ k=1,\ldots,N,
\]

we have that each of the functions $w_k$ satisfies (17.36) so that by summation over $k\neq l$ for some fixed $l$, we obtain

\[
\left\{R^{-n}\int_{B_R}\left[\sum_{k\neq l}(M_{2k}-w_k)\right]^{p}\right\}^{1/p}
\leq N^{1/p}\sum_{k\neq l}\left\{R^{-n}\int_{B_R}(M_{2k}-w_k)^{p}\right\}^{1/p}
\]
\[
\leq c\left\{\sum_{k\neq l}(M_{2k}-M_{1k})+R^2|D^2g|_{0;\Omega}\right\}
\]
\[
\leq C\left\{\omega(2R)-\omega(R)+R^2|D^2g|_{0;\Omega}\right\}
\]

where, for $s=1,2$,

\[
\omega(sR)=\sum_{k=1}^{N}\osc_{B_{sR}} w_k=\sum_{k=1}^{N}(M_{sk}-m_{sk})
\]

By (17.39), we have for $x\in B_{2R},\ y\in B_R$

\[
\beta_l\bigl(w_l(y)-w_l(x)\bigr)\leq g(y)-g(x)+\sum_{k\neq l}\beta_k\bigl(w_k(x)-w_k(y)\bigr)
\]


% PDF page 468
\source{gilbarg-trudinger:page-0468}
\providecommand{\osc}{\operatorname{osc}}
\providecommand{\D}{D}
\providecommand{\R}{\mathbb{R}}
\providecommand{\I}{I}
\providecommand{\g}{g}
\providecommand{\w}{\omega}
\providecommand{\Br}[1]{B_{#1}}

so that

\[
w_l(y)-m_{2k}\le \frac{1}{\lambda^*}\left\{3R|Dg|_{0;\Omega}+\Lambda^*\sum_{k\ne l}(M_{2k}-w_k)\right\};
\]

and hence

\begin{equation}\tag{17.40}
\left\{R^{-n}\int_{B_R}(w_l-m_{2l})^p\right\}^{1/p}
\le C\{ \omega(2R)-\omega(R)+R|Dg|_{0;\Omega}+R^2|D^2g|_{2;\Omega}\},
\end{equation}

where $C$ again depends only on $n$, $\lambda$, $\Lambda$. By setting $w=w_l$ in (17.36), adding this to (17.40) and summing over $l=1,\ldots,N$, we therefore obtain

\[
\omega(2R)\le C\{\omega(2R)-\omega(R)+R|Dg|_{0;\Omega}+R^2|D^2g|_{0;\Omega}\},
\]

hence

\[
\omega(R)\le \delta\omega(2R)+R|Dg|_{0;\Omega}+R^2|D^2g|_{0;\Omega}
\]

for $\delta=1-1/C$. H\"older estimates for the functions $w_k$, $k=1,\ldots,N$ now follow from Lemma 8.23, and by using the last assertion of Lemma 17.13, we obtain a H\"older estimate for $D^2u$: For any ball $B_{R_0}\subset\Omega$ and $R\le R_0$,

\begin{equation}\tag{17.41}
\osc_{B_R}D^2u \le C\left(\frac{R}{R_0}\right)^\alpha
\left\{\osc_{B_{R_0}}D^2u+R_0|Dg|_{0;\Omega}+R_0^2|D^2g|_{0;\Omega}\right\}
\end{equation}

where $C$ and $\alpha$ are positive constants depending only on $n$, $\lambda$ and $\Lambda$.

In terms of the interior H\"older norms the estimate (17.41) may be expressed in the form

\begin{equation}\tag{17.42}
|u|_{2,\alpha;\Omega}^* \le C\left(|u|_{2;\Omega}^*+|g|_{2;\Omega}^*\right)
\end{equation}

where $C$, $\alpha$ depend on $n$, $\lambda$ and $\Lambda$.

Let us now proceed to the general case (17.1) with $F\in C^2(I)$. Corresponding to conditions (i) and (ii) we suppose:

\begin{enumerate}
\item[(i)$'$] $F$ is uniformly elliptic with respect to $u$, so that there exist positive constants $\lambda$, $\Lambda$ such that
\end{enumerate}

\begin{equation}\tag{17.43}
\lambda|\xi|^2\le F_{ij}(x,u,Du,D^2u)\,\xi_i\xi_j\le \Lambda|\xi|^2
\end{equation}

for all $\xi\in\R^n$;

\begin{enumerate}
\item[(ii)] $F$ is concave with respect to $r$ on the range of $D^2u$.
\end{enumerate}

% PDF page 469
\source{gilbarg-trudinger:page-0469}


Again we differentiate twice with respect to a unit vector $\gamma$, to obtain the equations

\begin{align}
F_{ij}D_{ij\gamma}u + F_{p_i}D_{i\gamma}u + F_zD_{\gamma}u + \gamma_iF_{x_i} &= 0, \\
F_{ij}D_{ij\gamma\gamma}u + F_{ij,kl}D_{ij\gamma}uD_{kl\gamma}u + 2F_{ij,p_k}D_{ij\gamma}uD_{k\gamma}u \notag \\
\qquad + 2F_{ij,z}D_{ij\gamma}uD_{\gamma}u + 2\gamma_kF_{ij,x_k}D_{ij\gamma}u + F_{p_i}D_{i\gamma\gamma}u \notag \\
\qquad + F_{p_ip_j}D_{i\gamma}uD_{j\gamma}u + 2F_{p_i z}D_{i\gamma}uD_{\gamma}u + 2\gamma_jF_{p_i,x_j}D_{i\gamma}u \notag \\
\qquad + F_zD_{\gamma\gamma}u + F_{zz}(D_{\gamma}u)^2 + 2\gamma_iF_{z x_i}D_{\gamma}u + \gamma_i\gamma_jF_{x_i x_j} &= 0.
\end{align}

Using the concavity of $F$ we now obtain, in place of (17.35), the differential inequality

\begin{equation}
F_{ij}D_{ij\gamma\gamma}u \geq -A_{ij}D_{ij\gamma}u - B_{\gamma},
\end{equation}

where

\begin{align*}
A_{ij} &= 2F_{ij,p_k}D_{k\gamma}u + 2F_{ij,z}D_{\gamma}u + 2\gamma_kF_{ij,x_k} + \gamma_jF_{p_i}, \\
B_{\gamma} &= F_{p_ip_j}D_{i\gamma}uD_{j\gamma}u + 2F_{p_i z}D_{i\gamma}uD_{\gamma}u + 2\gamma_jF_{p_i,x_j}D_{i\gamma}u \\
&\qquad + F_zD_{\gamma\gamma}u + F_{zz}(D_{\gamma}u)^2 + 2\gamma_iF_{z x_i}D_{\gamma}u + \gamma_i\gamma_jF_{x_i x_j}.
\end{align*}

The third derivatives of $u$ in (17.45) are handled analogously to the second derivatives of $u$ in the derivation of the Hölder gradient estimate (Theorem 13.6). Let us first choose directions $\gamma_1,\dots,\gamma_N$ in accordance with the full statement of Lemma 17.13 applied to the matrix $[F_{ij}]$. Set

\[
M_2 = \sup_{\Omega} \lvert D^2u\rvert,
\]

\[
h_k = \frac12\left(1 + \frac{D_{\gamma_k\gamma_k}u}{1+M_2}\right), \qquad k=1,\dots,N,
\]

(if necessary we replace $\Omega$ by a subdomain to ensure the finiteness of $M_2$). We obtain from (17.45)

\begin{equation}
-F_{ij}D_{ij}h_k \leq C(A_0\lvert D^3u\rvert + B_0)/(1 + M_2)
\end{equation}

where $C=C(n)$ and

\begin{align*}
A_0 &= \sup_{\Omega}\left\{\lvert F_{rp}\rvert \lvert D^2u\rvert + \lvert F_{rz}\rvert \lvert Du\rvert + \lvert F_{rx}\rvert + \lvert F_p\rvert\right\}, \\
B_0 &= \sup_{\Omega}\left\{\lvert F_{pp}\rvert \lvert D^2u\rvert^2 + \lvert F_{pz}\rvert \lvert D^2u\rvert \lvert Du\rvert + \lvert F_{px}\rvert \lvert D^2u\rvert \right. \\
&\qquad\left. + \lvert F_z\rvert \lvert D^2u\rvert + \lvert F_{zz}\rvert \lvert Du\rvert^2 + \lvert F_{zx}\rvert \lvert Du\rvert + \lvert F_{xx}\rvert\right\},
\end{align*}

where $F_{rp} = [F_{ij,p_k}]_{i,j,k=1,\dots,n}$, etc., are evaluated at $(x,u,Du,D^2u)$. We now multiply (17.46) by $h_k$ and sum from $1$ to $N$ to obtain

\begin{equation}
\sum_{k=1}^N F_{ij}D_i h_k D_j h_k - \frac12 F_{ij}D_{ij}v \leq C(A_0\lvert D^3u\rvert + B_0)/(1 + M_2),
\end{equation}


% PDF page 470
\source{gilbarg-trudinger:page-0470}

where
\[
v=\sum_{k=1}^N (h_k)^2.
\]

By our choice of $\gamma_k$, we can estimate
\[
|D^3u|^2=\sum_{i,j,l=1}^n |D_{ijl}u|^2
\le 4(1+M_2)^2 \sum_{k=1}^N |Dh_k|^2,
\]
and by the ellipticity condition (17.43),
\[
\sum_{k=1}^N F_{ij}D_i h_k\,D_j h_k \geq \lambda \sum_{k=1}^N |Dh_k|^2.
\]

Consequently, for $\varepsilon\in(0,1)$ and
\[
w=w_k=h_k+\varepsilon v,\qquad k=1,\dots,N,
\]
we have, by combining (17.46) and (17.47),
\[
\varepsilon\lambda\sum_{k=1}^N |Dh_k|^2-\frac12 F_{ij}D_{ij}w
\le C\left\{A_0\left(\sum_{k=1}^N |Dh_k|^2\right)^{1/2}+\frac{B_0}{1+M_2}\right\};
\]
and hence, by the Cauchy inequality,
\begin{equation}
\tag{17.48}
F_{ij}D_{ij}w\geq -\lambda \bar{\mu},
\end{equation}
where
\[
\bar{\mu}=\frac{C(n)}{\lambda}\left(\frac{A_0^2}{\lambda\varepsilon}+\frac{B_0}{1+M_2}\right).
\]

We are now ready to apply again the weak Harnack inequality (Theorem 9.22).
Let $B_R$, $B_{2R}$ be concentric balls in $\Omega'$ and set, for $s=1,2$, $k=1,\dots,N$,
\[
W_{ks}=\sup_{B_{sR}} w,
\]
\[
M_{ks}=\sup_{B_{sR}} h_k,
\]
\[
m_{ks}=\inf_{B_{sR}} h_k,
\]
\[
\omega(sR)=\sum_{k=1}^N \osc_{B_{sR}} h_k=\sum_{k=1}^N (M_{ks}-m_{ks}).
\]


% PDF page 471
\source{gilbarg-trudinger:page-0471}

Applying Theorem 9.22 to the functions $W_{k2}-w_{k}$, we thus obtain

\begin{equation}
(17.49)\qquad \Phi_{p,R}(W_{k2}-w_{k})\equiv \left\{\frac{1}{|B_{R}|}\int_{B_{R}}(W_{k2}-w_{k})^{p}\right\}^{1/p}\le C\{W_{k2}-W_{k1}+\mu R^{2}\},
\end{equation}

where $p$ and $C$ are positive constants depending only on $n$ and $\Lambda/\lambda$. Using the inequalities

\[
W_{k2}-w_{k}\ge M_{k2}-h_{k}-2\epsilon\omega(2R),
\]

\[
W_{k2}-w_{k1}\le M_{k2}-M_{k1}+2\epsilon\omega(2R),
\]

we can conclude from (17.49) a similar inequality for the functions $h_{k}$; namely

\[
\Phi_{p,R}(M_{k2}-h_{k})\le C\{M_{k2}-M_{k1}+\epsilon\omega(2R)+\mu R^{2}\}.
\]

Then summation over $k\neq l$ for some fixed $l$ yields

\begin{equation}
(17.50)\qquad \Phi_{p,r}\left(\sum_{k\neq l}(M_{k2}-h_{k})\right)\le N^{1/p}\sum_{k\neq l}(M_{k2}-h_{k})
\le C\{(1+\epsilon)\omega(2R)-\omega(R)+\mu R^{2}\},
\end{equation}

where $C=C(n,\Lambda/\lambda)$ as before. To compensate for not having the inequality corresponding to (17.48) for the functions $-h_{k}$, we again resort to equation (17.1), so that using the concavity of $F$ (condition (ii)'), we have

\[
F_{ij}(y,u(y),Du(y),D^{2}u(y))(D_{ij}u(y)-D_{ij}u(x))
\]
\[
\le F(y,u(y),Du(y),D^{2}u(x)) - F(y,u(y),Du(y),D^{2}u(y))
\]
\[
= F(y,u(y),Du(y),D^{2}u(x)) - F(x,u(x),Du(x),D^{2}u(x))
\]
\[
\le D_{0}|x-y|,
\]

where

\[
D_{0}=\sup_{x,y\in\Omega}\Bigl\{|F_{x}(y,u(y),Du(y),D^{2}u(x))|
+|F_{z}(y,u(y),Du(y),D^{2}u(x))||Du(y)|
+|F_{p}(y,u(y),Du(y),D^{2}u(x))||D^{2}u(y)|\Bigr\}
\]

Now by Lemma 17.13 and our choice of $\gamma_{k}$,

\[
F_{ij}(y,u(y),Du(y),D^{2}u(y))(D_{ij}u(y)-D_{ij}u(x))
\]
\[
= \sum_{k=1}^{N}\beta_{k}(y)\bigl(D_{\gamma_{k}\gamma_{k}}u(y)-D_{\gamma_{k}\gamma_{k}}u(x)\bigr)
\]

\[
= 2(1+M_{2})\sum_{k=1}^{N}\beta_{k}(h_{k}(y)-h_{k}(x)),
\]

% PDF page 472
\source{gilbarg-trudinger:page-0472}

where

\[
0 < \lambda^* \le \beta_k \le \Lambda^*, \qquad k = 1,\dotsc,N,
\]

and $\lambda^*/\lambda$, $\Lambda^*/\lambda$ depend only on $n$ and $\Lambda/\lambda$. Consequently, for $x \in B_{2R}$, $y \in B_R$,

\[
\sum_{k=1}^N \beta_k\bigl(h_k(y) - h_k(x)\bigr) \le C \lambda \tilde{\mu} R,
\]

where

\[
\tilde{\mu} = \frac{D_0}{\lambda(1+M_2)};
\]

and hence for fixed $l$,

\[
h_l(y) - m_{12} \le \frac{1}{\lambda^*}\left\{C\lambda \tilde{\mu} R + \Lambda^*\sum_{k \ne l}\bigl(M_{k2} - h_k(y)\bigr)\right\}
\]

\[
\le C\left\{\tilde{\mu}R + \sum_{k \ne l}\bigl(M_{k2} - h_k(y)\bigr)\right\},
\]

where $C = C(n,\Lambda/\lambda)$. Therefore by (17.50), we obtain, for $l = 1,\dotsc,N$,

\[
\phi_{p,R}(h_l - m_{12}) \le C\bigl\{(1+\varepsilon)\omega(2R) - \omega(R) + \tilde{\mu}R + \tilde{\mu}R^2\bigr\},
\]

where $C = C(n,\Lambda/\lambda)$. By adding the above inequality for $l = k$ to (17.49) and summing over $k$, we thus obtain

\[
\omega(2R) \le C\bigl\{(1+\varepsilon)\omega(2R) - \omega(R) + \tilde{\mu}R + \tilde{\mu}R^2\bigr\};
\]

hence

\[
\omega(R) \le \delta\omega(2R) + C\bigl\{\varepsilon\omega(2R) + \tilde{\mu}R + \tilde{\mu}R^2\bigr\}
\]

for $\delta = 1 - 1/C$. Finally, by fixing $\varepsilon = (1-\delta)/2C$ we arrive again at the oscillation estimate

\[
\omega(R) \le \bar{\delta}\,\omega(2R) + C\bigl(\tilde{\mu}R + \tilde{\mu}R^2\bigr),
\]

where $0 < \bar{\delta} < 1$ and $C,\bar{\delta}$ depend only on $n$ and $\Lambda/\lambda$. The desired Hölder estimates now follow immediately from Lemma 8.23. Namely, for any ball $B_{R_0} \subset \Omega$ and $R \le R_0$, we have

\[
\tag{17.51}
\osc_{B_R} D^2u \le C\left(\frac{R}{R_0}\right)^\alpha (1+M_2)(1+\tilde{\mu}R_0+\tilde{\mu}R_0^2)
\]

% PDF page 473
\source{gilbarg-trudinger:page-0473}

\providecommand{\R}{\mathbb{R}}
\providecommand{\diam}{\operatorname{diam}}
\providecommand{\dist}{\operatorname{dist}}
\providecommand{\abs}[1]{\left|#1\right|}

where $C$ and $\alpha$ are positive constants depending only on $n$ and $\Lambda/\lambda$. The resultant interior estimate is formulated in the following theorem.

\begin{quote}
\textbf{Theorem 17.14.} \emph{Let $u \in C^4(\Omega)$ satisfy $F[u] = 0$ in $\Omega$ where $F \in C^2(\Gamma)$, $F$ is elliptic with respect to $u$ and satisfies $(i)'$ and $(ii)'$. Then for any $\Omega' \Subset \Omega$, we have the estimate}
\end{quote}

\begin{equation}\tag{17.52}
\bigl[D^2u\bigr]_{\alpha;\Omega'} \le C
\end{equation}

\noindent where $\alpha$ depends only on $n$, $\lambda$ and $\Lambda$, and $C$ depends in addition on $\abs{u}_{2;\Omega}$, $\dist(\Omega', \partial\Omega)$ and the first and second derivatives of $F$ other than $F_{rr}$.

A more explicit form of the estimate (17.52) is provided by (17.51), which exhibits the dependence of the constant $C$ on $\abs{u}_{2;\Omega}$ and the derivatives of $F$. Under further concavity assumptions on $F$ various terms in the expressions for $A_0$ and $B_0$ in (17.46) may be removed. For example, when $F$ is jointly concave in $z$, $p$ and $r$, we may take in (17.45)

\begin{align*}
A_{ij} &= 2\gamma_k F_{ij,x_k} + \gamma_j F_{p_i}, \\
B_\gamma &= 2\gamma_j F_{p_i,x_j} D_{iy}u + F_z D_{\gamma\gamma}u + 2\gamma_i F_{z,x_i} D_\gamma u + \gamma_i\gamma_j F_{x_i x_j}
\end{align*}

\noindent and hence in (17.46)

\begin{align*}
A_0 &= \sup_{\Omega}\bigl\{\abs{F_{rx}} + \abs{F_p}\bigr\}, \\
B_0 &= \sup_{\Omega}\bigl\{\abs{F_{px}}\abs{D^2u} + \abs{F_z}\abs{D^2u} + \abs{F_{zx}}\abs{Du} + \abs{F_{xx}}\bigr\}.
\end{align*}

Consequently, under the additional structure conditions

\begin{equation}\tag{17.53}
0 < \lambda \abs{\xi}^2 \le F_{ij}(x,z,p,r)\,\xi_i \xi_j \le \Lambda \abs{\xi}^2
\end{equation}
\begin{equation*}
\abs{F_p}, \abs{F_z}, \abs{F_{rx}}, \abs{F_{px}}, \abs{F_{zx}} \le \mu \lambda,
\end{equation*}
\begin{equation*}
\abs{F_x}, \abs{F_{xx}} \le \mu\lambda(1 + \abs{p} + \abs{r})
\end{equation*}

\noindent for all nonzero $\xi \in \R^n$, $(x,z,p,r) \in \Gamma$, where $\lambda$ is a nonincreasing function of $\abs{z}$ and $\Lambda$ and $\mu$ are nondecreasing functions of $\abs{z}$, we may conclude from (17.51) the interior estimate, extending (17.42),

\begin{equation*}
\abs{u}^{*}_{2,\alpha} \le C(1 + \abs{u}^{*}_{2}),
\end{equation*}

\noindent where $\alpha$ depends only on $n$, $\Lambda/\lambda$, and $C$ depends in addition on $\mu$, $\diam \Omega$. With the aid of the interior interpolation inequality (Lemma 6.32), we thus obtain the following interior estimate.

% PDF page 474
\source{gilbarg-trudinger:page-0474}

\providecommand{\R}{\mathbb{R}}
\providecommand{\abs}[1]{\left|#1\right|}
\providecommand{\diam}{\operatorname{diam}}

\textbf{Theorem 17.15.} \emph{Let $u \in C^4(\Omega)$ satisfy $F[u] = 0$ in $\Omega \subset \R^n$ and suppose that $F$ is concave (or convex) in $z$, $p$, $r$ and satisfies the structure conditions (17.53). Then we have the interior estimate}

\begin{equation}\tag{17.54}
\abs{u}^{*}_{2,\alpha;\Omega} \le C
\end{equation}

\noindent \emph{where $\alpha > 0$ depends only on $n$ and $\Lambda(M)/\lambda(M)$, and $C$ depends in addition on $\mu(M)$, $\diam \Omega$ and $M = \abs{u}_{0;\Omega}$.}

\medskip

We remark that the estimate (17.54) can in fact be proved under more general hypotheses corresponding to the natural conditions for quasilinear equations; (see Notes). We conclude this section with the proof of Lemma 17.13.

\medskip

\textit{Proof of Lemma 17.13.} \quad Letting $\R^{n\times n}_+$ denote the cone of positive matrices in $\R^{n\times n}$, we may represent any $\mathscr{A} \in \R^{n\times n}_+$ in the form

\begin{equation}\tag{17.55}
\mathscr{A} = \sum_{k=1}^{n'} \gamma_k \otimes \gamma_k
\end{equation}

\noindent where $n' = n(n + 1)/2 = \dim.\ \R^{n\times n}$, $\gamma_k \in \R^n$, $k = 1, . . . , n'$, and the dyadic matrices
$\gamma_k \otimes \gamma_k = [\gamma_{ki}\gamma_{kj}]$ are linearly independent. To see this we observe that any two matrices in $\R^{n\times n}_+$ are similar and hence in particular each $\mathscr{A} \in \R^{n\times n}_+$ is similar to the matrix $\mathscr{A}_0$ whose diagonal and nondiagonal terms are, respectively, $n$ and $1$. But then

\begin{equation*}
\mathscr{A}_0 = \sum_{i=1}^n e_i \otimes e_i + \sum_{\substack{i,j=1 \\ i<j}}^n (e_i + e_j) \otimes (e_i + e_j)
\end{equation*}

\noindent so that (17.55) follows by an appropriate base change. Consequently, the family of sets of the form

\begin{equation*}
U(\gamma_1, \dots, \gamma_{n'}) = \left\{ \sum_{k=1}^{n'} \beta_k \gamma_k \otimes \gamma_k \,\middle|\, \beta_k > 0,\, k = 1, . . . , n' \right\},
\end{equation*}

\noindent where $\gamma_k \otimes \gamma_k$ are linearly independent, forms an open cover of $S(\lambda, \Lambda) \subset \R^{n\times n}_+$, and since $S(\lambda, \Lambda)$ is compact there exists a finite subcover. Accordingly, there exists a fixed set of unit vectors $\gamma_1, \dots, \gamma_N$, depending only on $\lambda$, $\Lambda$ and $n$ such that any $\mathscr{A} \in S(\lambda, \Lambda)$ may be written

\begin{equation*}
\mathscr{A} = \sum_{k=1}^N \beta_k \gamma_k \otimes \gamma_k
\end{equation*}

\noindent with $\beta_k \ge 0$. To get the assertion of the lemma we observe that at the outset we could have applied the above procedure to the matrix

\begin{equation*}
\mathscr{A} - \lambda^{*} \sum_{k=1}^N \gamma_k \otimes \gamma_k \in S(\lambda/2, \Lambda)
\end{equation*}

% PDF page 475
\source{gilbarg-trudinger:page-0475}
\begin{flushleft}
for sufficiently small $\lambda^*$ ($\lambda^*=\lambda/2N$ is sufficient). Note that we can take $A^*=A$. A similar consideration shows that any particular finite set of unit vectors may be included among the $\gamma_k$. \qed
\end{flushleft}

\section*{17.5. Dirichlet Problem for Uniformly Elliptic Equations}

We show in this section that the interior derivative estimates of the preceding section in fact suffice to demonstrate the solvability of the Dirichlet problem for certain types of uniformly elliptic equations, including those of Bellman-Pucci type. Since these estimates were established for $C^4$ solutions, we first need a regularity result to link them with the hypotheses of the method of continuity as asserted in Theorem 17.8.

\paragraph{Lemma 17.16.}
Let $u \in C^2(\Omega)$ satisfy $F[u]=0$ in $\Omega$ where $F$ is elliptic with respect to $u$. Then, if $F \in C^k(\Gamma)$, $k \ge 1$, we have $u \in W^{k+2,p}_{\loc}(\Omega)$ for all $p<\infty$; if $F \in C^{k,\alpha}(\Gamma)$, $0<\alpha<1$, we have $u \in C^{k+2,\alpha}(\Omega)$.

\paragraph{Proof.}
We use a difference quotient argument similar to that in the proof of Theorem 6.17. Let us fix a coordinate vector $e_\ell$, $1 \le \ell \le n$, and write

\[
v(x)=u(x+he_\ell), \qquad h \in \R
\]
\[
w=\Delta_\ell^h u=\frac{1}{h}(u-v),
\]
\[
u_\theta=\theta u+(1-\theta)v, \qquad 0 \le \theta \le 1,
\]
\[
a^{ij}(x)=\int_0^1 F_{ij}(x+\theta h,u_\theta,Du_\theta,D^2u_\theta)\,d\theta,
\]
\[
b^i(x)=\int_0^1 F_{p_i}(x+\theta h,u_\theta,Du_\theta,D^2u_\theta)\,d\theta,
\]
\[
c(x)=\int_0^1 F_z(x+\theta h,u_\theta,Du_\theta,D^2u_\theta)\,d\theta,
\]
\[
f(x)=\int_0^1 F_{x_\ell}(x+\theta h,u_\theta,Du_\theta,D^2u_\theta)\,d\theta.
\]

Then if we fix a subdomain $\Omega' \Subset \Omega$ and take sufficiently small $h$, the difference quotient $w$ will satisfy in $\Omega'$ the linear equation

\begin{equation}
Lw=a^{ij}D_{ij}w+b^iD_iw+cw=-f,
\tag{17.56}
\end{equation}

% PDF page 476
\source{gilbarg-trudinger:page-0476}
\begin{flushleft}
which will also be elliptic with uniformly continuous coefficients in $\Omega'$. The interior
$L^p$ estimates (Theorem 9.11) then yield bounds for $\|D^2v\|_{p;\Omega'}$ for $\Omega'' \subset \subset \Omega'$,
independent of $h$, and hence by Lemma 7.24 we conclude $u \in W^{3,p}_{\loc}(\Omega)$ for any
$p<\infty$. By the Sobolev imbedding theorem (Theorem 7.10), it then follows that
$u \in C^{2,\alpha}(\Omega)$ for all $\alpha<1$. Consequently, if $F \in C^{1,\alpha}(\Gamma)$ for some $\alpha$, $0<\alpha<1$, we
may apply the Schauder regularity result (Theorem 6.17) to (17.56) to obtain
$u \in C^{3,\alpha}(\Omega)$. Note that if at the outset we are given $u \in C^{2,\beta}(\Omega)$, for some $\beta>0$, there
is no need to use the $L^p$ theory for this result. Lemma 17.16 is thus established for
$k=1$. Further interior regularity now follows by a straightforward iterative or
``bootstrapping'' procedure. \qed
\end{flushleft}

\begin{flushleft}
With the aid of Lemma 17.16 we see that the estimates of Sections 17.3 and 17.4
will hold for $C^2(\Omega)$ solutions of (17.1). In the case of Section 17.4, we should observe
that $u \in W^{4,n}_{\loc}(\Omega)$ is sufficient for the given proofs. To offset the lack of global $C^{2,\alpha}(\bar{\Omega})$
estimates, we modify the function $F$ near the boundary $\partial\Omega$ as follows. Let $\{\eta_m\}$ be a
sequence of functions in $C^2_0(\Omega)$ satisfying $0 \leq \eta \leq 1$ in $\Omega$ and $\eta_m(x)=1$ whenever
$d(x,\partial\Omega)\geq 1/m$, and consider instead of $F$ the operators $F_m$ given by
\end{flushleft}

\[
(17.57)\qquad F_m[u]=\eta_m F[u]+(1-\eta_m)\Delta u.
\]

\begin{flushleft}
If $F$ satisfies the hypotheses of Theorems 17.14 or 17.15, then so also does $F_m$ with
structural constants depending possibly on $\eta_m$. We therefore obtain interior $C^{2,\alpha}$
estimates for solutions of the equations, $F_m[u]=0$, of the same form as those given
by Theorems 17.14 and 17.15. But near $\partial\Omega$, $F_m[u]=\Delta u$ so that for appropriately
smooth boundary data, $C^{2,\alpha}$ estimates near $\partial\Omega$ will follow from the Schauder theory,
in particular from Lemma 6.5. Using this procedure we can now establish the follow-
ing existence result.
\end{flushleft}

\paragraph{Theorem 17.17.}
Let $\Omega$ be a bounded domain in $\R^n$ satisfying an exterior sphere con-
dition at each boundary point and suppose the function $F \in C^2(\Gamma)$ is concave (or con-
vex) with respect to $z$, $p$, $r$, nonincreasing with respect to $z$ and satisfies the structure
conditions (17.53). Then the classical Dirichlet problem, $F[u]=0$ in $\Omega$, $u=\phi$ on
$\partial\Omega$ is uniquely solvable in $C^2(\Omega)\cap C^0(\bar\Omega)$ for any $\phi \in C^0(\partial\Omega)$.

\paragraph{Proof.}
Let us first consider the case of smooth boundary data, namely $\partial\Omega \in C^{2,\beta}$,
$\phi \in C^{2,\beta}(\bar\Omega)$ for some $0<\beta<1$. In view of the discussion preceding the theorem,
we consider the approximating Dirichlet problems,

\[
(17.58)\qquad F_m[u]=0,\qquad u=\phi\quad \text{on }\partial\Omega
\]

\begin{flushleft}
where $F_m$ is given by (17.57). In order to apply the method of continuity (Theorem
17.8), we require apriori bounds in $C^{2,\alpha}(\bar\Omega)$ for some $\alpha>0$ for solutions of the
problems
\end{flushleft}

\[
(17.59)\qquad F_m[u]-tF_m[\phi]=0,\qquad u=\phi\quad \text{on }\partial\Omega,\qquad 0\leq t\leq 1.
\]

% PDF page 477
\source{gilbarg-trudinger:page-0477}

Now the equations (17.59) continue to satisfy the hypotheses of Theorem 17.14,
uniformly with respect to $t$; and consequently we have, for any solution $u \in C^{2,\beta}(\bar\Omega)$
and $\Omega' \Subset \Omega$, an estimate
\[
|u|_{2,\alpha;\Omega'} \leq C,
\]
where $\alpha$ depends only on $n$, $\lambda(M)$, $\Lambda(M)$, and $C$ depends in addition on $\mu(M)$, $\Omega$, $\Omega'$, $\phi$, $\eta_m$ and $M = |u|_{0;\Omega}$. Since $F_m[u] = \Delta u$ near $\partial\Omega$, we therefore conclude, by Lemma 6.5, the corresponding global estimate with $\Omega' = \Omega$ and $C$ depending in addition on $\partial\Omega$ and $\beta$. Next we observe that the conditions, $F_z \leq 0$ and (17.53) together imply (17.11), for $\mu_1 = \mu_2 = \mu(0)$, so that by Theorem 17.3, $M = |u|_{0;\Omega}$ is bounded uniformly in $t$ and $m$. Therefore, by Theorem 17.8, we obtain the existence of a unique solution $u = u_m \in C^{2,\beta}(\bar\Omega)$ of the Dirichlet problem (17.58). Since $F_m[u] = F[u]$ for $\dist(x, \partial\Omega) \geq 1/m$, we then obtain, by the interior estimate (Theorem 17.15), the convergence of a subsequence of $\{u_m\}$ (uniformly together with first and second derivatives on compact subsets of $\Omega$) to a solution $u \in C^{2,\beta}(\Omega)$ of the equation $F[u] = 0$ in $\Omega$. But in view of the representation (17.10) the barrier considerations for quasilinear equations, in particular Theorem 14.15, are applicable to the Dirichlet problems (17.58) and as a result we obtain that $u \in C^0_1(\Omega)$ and $u = \phi$ on $\partial\Omega$. The extension to continuous $\phi$ and domains $\Omega$ satisfying exterior sphere conditions follows similarly. $\square$

By approximation, the condition $F \in C^2(\Gamma)$ can be weakened in the hypotheses of Theorem 17.17, so that the derivatives appearing in the structure conditions (17.53) need exist only in the weak sense; (see Problem 17.4). Furthermore, by extending the above arguments slightly we can encompass the uniformly elliptic Bellman equations (17.8). In fact the necessary modifications can be illustrated more generally. Let $F^1, \dots, F^m$ be operators of the form (17.1) and let $G \in C^2(\mathbb{R}^m)$ be a concave function whose derivatives satisfy
\[
(17.60)\qquad 1 \leq \sum_{\nu=1}^m D_\nu G \leq K
\]
for some constant $K$. We now define another operator $F$ by
\[
F[u] = G(F^1[u], \dots, F^m[u]).
\]
Then if $u \in C^4(\Omega)$ and $F^1, \dots, F^m$ all satisfy hypotheses (i)' and (ii)' in Theorem 17.14, we obtain by differentiation, in place of (17.44), the equation
\[
\sum_{\nu=1}^m D_\nu G\, D_\nu F^\nu = 0,
\]
\[
\sum_{\nu=1}^m D_\nu G\, D_{\nu\nu} F^\nu + \sum_{\nu,\tau=1}^m D_{\nu\tau} G\, D_\nu F^\nu D_\tau F^\tau = 0,
\]

% PDF page 478
\source{gilbarg-trudinger:page-0478}

so that from the concavity of $G$ we have

\[
\sum_{\nu=1}^m D_\nu G D_{\gamma\gamma}F^\nu \ge 0.
\]

Consequently the analysis of Section 17.4 carries over with the derivatives $DF,\, D^2F$,
replaced respectively by $\sum_{\nu=1}^m D_\nu GDF^\nu$, $\sum_{\nu=1}^m D_\nu GD^2F^\nu$. Using (17.60) we thus have, in
particular, that the estimate of Theorem 17.14 is applicable to the operator $F$
provided the operators $F^\nu$, $\nu = 1,\dots,m$, satisfy the structural hypotheses and $\lambda$
and $\mu$ are replaced by $K\lambda$ and $K\mu$ respectively. The existence result (Theorem 17.17)
extends similarly. The equations of Bellman-Pucci type can then be treated by
approximation. Namely, let us define for $y \in \R^m$,

\[
G_0(y)=\inf_{\nu=1,\dots,m} y_\nu
\]

and for $h>0$, let $G_h$ be the mollification of $G_0$ given by

\[
G_h(y)=h^{-n}\int_{\R^m}\rho\!\left(\frac{y-\bar y}{h}\right)G_0(\bar y)\,d\bar y,
\]

where $\rho$ is a mollifier on $\R^m$. Since $G_0$ is concave, it follows readily that $G_h$ is also
concave. Furthermore we have

\[
\sum_{\nu=1}^m D_\nu G_h = 1,
\]

so that (17.60) holds with $K=1$. It follows then that if $F^1,\dots,F^m$, $\Omega$ and $\phi$ satisfy
the hypotheses of Theorem 17.17, the classical Dirichlet problems,

\[
F[u]=G_h(F^1[u],\dots,F^m[u])=0,\qquad u=\phi \quad \text{on } \partial\Omega
\]

are uniquely solvable with solutions $u=u_h$ satisfying an estimate

\[
\normstar{u}\le C
\]

where $\alpha$ and $C$ depend only on $n,\lambda,\Lambda,\mu,\phi$ and $\Omega$, together with a modulus of con-
tinuity estimate on $\partial\Omega$ depending on the same quantities. By approximation, the
result extends to the limiting case $h=0$ and, since the above estimates are inde-
pendent of $m$, also to the case of a countable family of operators. Accordingly we
have the following extension of Theorem 17.17.

\textbf{Theorem 17.18.} \textit{Let $\Omega$ be a bounded domain in $\R^n$ satisfying an exterior sphere
condition at each boundary point and suppose the functions $F^1$, $F^2$, $\dots \in C^2(\Gamma)$, are}

% PDF page 479
\source{gilbarg-trudinger:page-0479}



concave with respect to $z,p,r$, nonincreasing with respect to $z$, and satisfy uniformly
the structure conditions (17.53). Then the classical Dirichlet problem,

\begin{equation}
F[u]=\inf \{F^1[u],F^2[u],\ldots\}=0 \quad \text{in }\Omega,\qquad u=\phi \quad \text{on }\partial\Omega
\tag{17.61}
\end{equation}

is uniquely solvable in $C^2(\Omega)\cap C^0(\overline{\Omega})$ for any $\phi\in C^0(\partial\Omega)$.

We observe that the Bellman equation (17.8) is included in Theorem 17.18
provided the index set $V$ is countable and the operators $L_\nu$ and $f_\nu$ satisfy the conditions:
$a_\nu^{ij},\, b_\nu^i,\, c_\nu,\, f_\nu \in C^2(\Omega)$ and

\begin{equation}
\lambda |\xi|^2 \le a_\nu^{ij}\xi_i\xi_j \le \Lambda |\xi|^2
\tag{17.62}
\end{equation}
\[
|a_\nu^{ij}|_{2;\Omega},\,|b_\nu^i|_{2;\Omega},\,|c_\nu|_{2;\Omega},\,|f_\nu|_{2;\Omega}\le \mu\lambda
\]
\[
c_\nu\le 0
\]

for all $\xi\in \mathbb{R}^n$, $\nu\in V$, where $\lambda$, $\Lambda$ and $\mu$ are positive constants. Moreover it is evident
that we can allow certain types of uncountable sets $V$, for example a separable
metric space on which the mappings: $\nu\to a_\nu^{ij}(x),b_\nu^i(x),c_\nu(x),f_\nu(x)$ are continuous for
each $x\in\Omega$. In particular the Pucci equations (17.6) are covered provided $f\in
C^2(\Omega)\cap L^\infty(\Omega)$.

\section*{17.6. Second Derivative Estimates for Equations of Monge-Amp\`ere Type}

In this section we turn our attention to equations of form

\begin{equation}
\det D^2u = f(x,u,Du).
\tag{17.63}
\end{equation}

As indicated earlier, equation (17.63) is elliptic only when the Hessian matrix $D^2u$
is positive (or negative) and it is therefore natural to confine our attention to
convex solutions $u$ and positive functions $f$. Writing equation (17.63) in the form

\begin{equation}
F(D^2u)=\log \det D^2u = g(x,u,Du),
\tag{17.64}
\end{equation}

where $g=\log f$, we then have by a calculation,

\begin{equation}
F_{ij}=u^{ij}
\tag{17.65}
\end{equation}
\[
F_{ij,kl}=-u^{ik}u^{jl}=-F_{ik}F_{jl}
\]

where $[u^{ij}]$ denotes the inverse of $D^2u$. Consequently the function $F$ is concave on
the cone of nonnegative matrices in $\mathbb{R}^{n\times n}$, and the equation (17.64) will be uniformly
elliptic on compact subsets of $\Omega$, with respect to any convex $C^2(\Omega)$ solution. If
$g\in C^2(\Omega\times\mathbb{R}\times\mathbb{R}^n)$ we can then apply the results of Section 17.4 to get interior
H\"older estimates for the second derivatives of solutions and subsequently, through
Lemma 17.16, higher order estimates when $g$ is appropriately smooth. We consider


% PDF page 480
\source{gilbarg-trudinger:page-0480}

now the question of interior and global bounds for the second derivatives of solutions, with approach reminiscent of the treatment of gradient estimates for non-uniformly elliptic equations in Chapter 15.

First let us note from equation (17.44) that any pure second derivative $D_{\gamma\gamma}u$, of a solution of (17.64), satisfies the equation

\begin{equation}
(17.66)\qquad F_{ij}D_{ij\gamma\gamma}u = F_{ik}F_{jl}D_{ij\gamma}uD_{kl\gamma}u + D_{\gamma\gamma}g.
\end{equation}

and since $u$ is convex we have also $D_{\gamma\gamma}u > 0$. To estimate $D_{\gamma\gamma}u$ from above, we take positive functions $\eta\in C^2(\Omega)$, $h\in C^2(\mathbb{R}^n)$ and set

\[
w=\eta h(Du)D_{\gamma\gamma}u
\]

so that

\begin{align*}
\frac{D_iw}{w} &= \frac{D_i\eta}{\eta} + (\log h)_{p_k}D_{ik}u + \frac{D_{i\gamma\gamma}u}{D_{\gamma\gamma}u},\\
\frac{D_{ij}w}{w} &= \frac{D_iwD_jw}{w^2} + \frac{D_{ij}\eta}{\eta} - \frac{D_i\eta D_j\eta}{\eta^2}\\
&\quad + (\log h)_{p_kp_l}D_{ik}uD_{jl}u + (\log h)_{p_k}D_{ijk}u\\
&\quad + \frac{D_{ij\gamma\gamma}u}{D_{\gamma\gamma}u} - \frac{D_{i\gamma\gamma}uD_{j\gamma\gamma}u}{(D_{\gamma\gamma}u)^2}.
\end{align*}

Consequently, using (17.66) we obtain

\begin{equation}
(17.67)\qquad (\eta h)^{-1}F_{ij}D_{ij}w \ge D_{\gamma\gamma}u\left\{\frac{F_{ij}D_{ij}\eta}{\eta} - \frac{F_{ij}D_i\eta D_j\eta}{\eta^2}\right.
\end{equation}
\[
\left.
+ (\log h)_{p_kp_l}F_{ij}D_{ik}uD_{jl}u + (\log h)_{p_k}F_{ij}D_{ijk}u\right\}
\]
\[
+ F_{ik}F_{jl}D_{ij\gamma}uD_{kl\gamma}u - \frac{1}{D_{\gamma\gamma}u}F_{ij}D_{i\gamma\gamma}uD_{j\gamma\gamma}u + D_{\gamma\gamma}g.
\]

An obvious candidate for the function $h$ is given by

\[
h(p)=e^{\beta|p|^2/2},\qquad \beta>0.
\]

For then we have

\[
(\log h)_{p_k}=\beta p_k,\qquad (\log h)_{p_kp_l}=\beta\delta_{kl},
\]

and hence

\[
(\log h)_{p_kp_l}F_{ij}D_{ik}uD_{jl}u = \beta F_{ij}D_{ik}uD_{jk}u = \beta\Delta u
\]


% PDF page 481
\source{gilbarg-trudinger:page-0481}
\begin{quote}
by (17.65). Next, if we assume the bounds
\end{quote}

\[
\tag{17.68}
|Dg(x,u,Du)|,\qquad |D^2g(x,u,Du)|\le \mu,
\]

\begin{quote}
we have
\end{quote}

\[
\begin{aligned}
D_{\gamma\gamma}u(\log h)_{p_k}F_{ij}D_{ijk}u + D_{\gamma\gamma}g
&= \beta D_k u D_{\gamma\gamma}u(g_{x_k} + g_z D_k u + g_{p_i}D_{ik}u) + g_{\gamma\gamma} + 2g_{\gamma z}D_\gamma u + 2g_{\gamma p_i}D_{i\gamma}u \\
&\quad + g_{zz}(D_\gamma u)^2 + 2g_{zp_i}D_\gamma uD_{i\gamma}u + g_{p_ip_j}D_{ik}uD_{jk}u \\
&\quad + g_z D_{\gamma\gamma}u + g_{p_i}D_{i\gamma\gamma}u \\
&\ge g_{p_i}\left(\frac{D_iw}{w} - \frac{D_i\eta}{\eta}\right)D_{\gamma\gamma}u - C\{1 + |D^2u|^2 + \beta(1 + |D^2u|)\},
\end{aligned}
\]

\begin{quote}
where $C$ depends on $\mu$, $\sup |Du|$.
\end{quote}

\begin{quote}
In order to handle the other terms in (17.67) we regard $w = w(x,y)$ as a function on $\Omega \times \partial B_1(0)$ and suppose that $w$ takes a maximum value at a point $y \in \Omega$ and direction $\gamma$. The derivative $D_{\gamma\gamma}u(y)$ will then be the maximum eigenvalue of the Hessian $D^2u(y)$ and by a rotation of coordinates we can assume that $D^2u(y)$ is in diagonal form with $\gamma$ a coordinate direction. For global estimates we take $\eta \equiv 1$ so that terms involving $\eta$ are not present in (17.67). It then follows from (17.65) that
\end{quote}

\[
\frac{1}{D_{\gamma\gamma}u}F_{ij}D_{ijyy}uD_{\gamma\gamma}u \le F_{ik}F_{jl}D_{ijyy}uD_{kly}u
\]

\begin{quote}
at the point $y$, so that by choosing $\beta$ sufficiently large we obtain an estimate for $D_{\gamma\gamma}u(y)$ in terms of $n$, $\mu$ and $|Du|_{0;\Omega}$.
\end{quote}

\begin{quote}
The interior case is more delicate since $\eta$ cannot be chosen as an arbitrary cutoff function (because of the term $F_{ij}D_{ij}\eta$). In this case we shall assume that $u$ vanishes continuously on $\partial\Omega$ and choose $\eta = -u$ so that $\eta > 0$ in $\Omega$ and
\end{quote}

\[
F_{ij}D_{ij}\eta = -F_{ij}D_{ij}u = -n.
\]

\begin{quote}
Furthermore since $Dw(y) = 0$, we have
\end{quote}

\[
\frac{F_{ij}D_i\eta D_j\eta}{\eta^2}
= \frac{\sum F_{ii}|D_i\eta|^2}{\eta^2}
= \frac{|D_\gamma u|^2}{u^2D_{\gamma\gamma}u} + \sum_{i\ne \gamma}F_{ii}\left(\frac{D_{i\gamma\gamma}u}{D_{\gamma\gamma}u} + D_i u D_{ii}u\right)^2
\]

\[
\le \frac{|D_\gamma u|^2}{u^2D_{\gamma\gamma}u} + \sum_{i\ne \gamma}F_{ii}\left(\frac{D_{i\gamma\gamma}u}{D_{\gamma\gamma}u}\right)^2 - \frac{2\beta|Du|^2}{u}
\]

% PDF page 482
\source{gilbarg-trudinger:page-0482}
\providecommand{\dist}{\operatorname{dist}}
\providecommand{\diam}{\operatorname{diam}}

at the point $y$. Now

\[
\frac{1}{D_{\gamma\gamma}u}
\left\{
\sum_{i\neq\gamma} F_{ii}(D_{i\gamma\gamma}u)^2
+
F_{ij}D_{i\gamma\gamma}u\,D_{j\gamma\gamma}u
\right\}
\]
\[
= \sum_{i\neq\gamma} F_{\gamma\gamma}F_{ii}(D_{i\gamma\gamma}u)^2
+\sum_{i=1}^n F_{\gamma\gamma}F_{ii}(D_{i\gamma\gamma}u)^2
\]
\[
\le \sum_{i,j=1}^n F_{ii}F_{jj}(D_{ij\gamma}u)^2
\]
\[
= F_{ik}F_{jl}D_{ij\gamma}u\,D_{kl\gamma}u.
\]

at $y$, by virtue of our choice of coordinates. Taking account of the above estimates,
in the differential inequality (17.67), we obtain

\[
D_{\gamma\gamma}u(y)\le C\left(1-\frac{1}{u(y)}\right)
\]

where $C=C(n,\mu,|Du|_{0;\Omega})$, and hence

\begin{equation}\tag{17.69}
\sup_{\Omega} w \le C,
\end{equation}

where $C=C(n,\mu,|u|_{1;\Omega})$. Finally, to estimate $\eta=-u$ from below, we have, by the
convexity of $u$,

\begin{equation}\tag{17.70}
\frac{u(x)}{\dist(x,\partial\Omega)}\le \frac{\inf_{\Omega} u}{\diam \Omega}
\end{equation}

for any $x\in\Omega$. We formulate the resultant second derivative estimates as follows.

\textbf{Theorem 17.19.} \emph{Let $u\in C^2(\Omega)$ be a convex solution of (17.63) in a domain $\Omega$ where
$f\in C^2(\Omega\times\mathbb{R}\times\mathbb{R}^n)$ is positive in $\Omega$ and $g=\log f$ satisfies (17.68). Then if $u\in C^2(\overline{\Omega})$,
we have}

\begin{equation}\tag{17.71}
\sup_{\Omega} |D^2u| \le C,
\end{equation}

where $C$ depends on $n,\mu,|u|_{1;\Omega}$ and $\sup_{\partial\Omega}|D^2u|$. If $u\in C^{0,1}(\overline{\Omega})$ and $u$ is constant on $\partial\Omega$,
\emph{then for any $\Omega' \Subset \Omega$, we have}

\begin{equation}\tag{17.72}
\sup_{\Omega'} |D^2u| \le \frac{C}{d_{\Omega'}},
\end{equation}

where $C$ depends on $n,\mu,|u|_{1;\Omega}$ and $\diam \Omega$, and $d_{\Omega'} = \dist(\Omega',\partial\Omega)$.

% PDF page 483
\source{gilbarg-trudinger:page-0483}

17.7.\ Dirichlet Problem for Equations of Monge-Amp\`ere Type

When the function $f^{1/n}$ is convex with respect to the $p$ variables, the estimate
(17.71) may be derived by simpler means and moreover the result extends to more
general functions $F$ and solutions that are not necessarily convex (see Problem
17.5).

An estimate for the second derivatives of solutions of equation (17.64) on the
boundary $\delOh$ follows readily from equation (17.26). If we assume that $\Oh$ is uniformly
convex, $\delOh\in C^3$, and $u=\phi$ on $\delOh$ where $\phi\in C^3(\overline{\Oh})$, then the first term on the right-
hand side of (17.26) becomes

\[
(17.73)\qquad 2F_{ij}D_i\!\left(\frac{\partial x_i}{\partial y_k}\right)D_{jl}u
= 2\delta_{il}D_i\!\left(\frac{\partial x_i}{\partial y_k}\right)
= 2D_i\!\left(\frac{\partial x_i}{\partial y_k}\right),
\]

by virtue of (17.65). An estimate for the derivatives $D_{v_k v_n}u(x_0)$ for $x_0\in \Oh$ then follows
from Theorem 14.4 or Corollary 14.5. The remaining second derivative $D_{v_n v_n}u(x_0)$ is
estimated directly from equation (17.64). For, with respect to a principal coordinate
system at $x_0$, we have (taking $\phi\equiv 0$)

\[
(17.74)\qquad \det D^2u = |D_nu|^{\,n-2}\prod_{i=1}^{n-1}\kappa_i\left\{D_nu\,D_{nn}u-\sum_{i=1}^{n-1}\frac{(D_{in}u)^2}{\kappa_i}\right\}
\]

where $\kappa_1,\ldots,\kappa_{n-1}$ are the principal curvatures at $x_0$, and using (17.70), we thus
infer a bound for $D_{nn}u(x_0)$. Therefore, we have proved the following global estimate
for second derivatives.

\textbf{Theorem 17.20.} \textit{Let $u\in C^3(\overline{\Oh})$ be a convex solution of equation (17.64) in $\Oh$, where
$f\in C^2(\overline{\Oh}\times \RR \times \RR^n)$ is positive and $\delOh\in C^3$ is uniformly convex. Then}

\[
(17.75)\qquad \sup_{\Oh} |D^2u| \le C
\]

\textit{where $C$ depends on $n$, $|u|_{1;\Oh}$, $f$, $\delOh$, and $u=0$ on $\delOh$.}

We remark here that the proof of Theorem 17.20 also extends to more general
functions $F$, with possibly nonconvex solutions $u$; (see Problems 17.6, and 17.7).

17.7.\ Dirichlet Problem for Equations of
Monge-Amp\`ere Type

The considerations of the previous section reduce the solvability of the classical
Dirichlet problem for equations of Monge-Amp\`ere type to the establishment of $C^1$
estimates. For equations in two variables, we can then proceed directly through the
method of continuity using the global H\"older estimates for second derivatives


% PDF page 484
\source{gilbarg-trudinger:page-0484}


(Theorem 17.10). For higher dimensions, there also exist procedures requiring only
interior second derivative estimates but these are more complicated than the method
employed for the uniformly elliptic case in Section 17.5; (see Notes). In the next
section, we shall treat the recently established global Hölder estimates for second
derivatives that will enable us to align the general case with the two-variable case.

Restrictions on the growth of the function $f$ in (17.63) arise through the con-
sideration of gradient estimates. For a convex function $u$ in a domain $\Omega$ we clearly
have

\begin{equation}
\tag{17.76}
\sup_{\Omega} |Du| = \sup_{\partial\Omega} |Du|
\end{equation}

so that the estimation of the gradient of convex solutions of (17.63) reduces to
estimation on the boundary only. As in the quasilinear case, such estimates are
readily obtained by means of barrier constructions. In fact, let us assume the follow-
ing structure condition,

\begin{equation}
\tag{17.77}
0 \le f(x,z,p) \le \mu(|z|) d_x^{\beta} |p|^{\gamma}
\end{equation}

for all $x$ in some neighbourhood $\mathcal N$ of $\partial\Omega$, $z \in \mathbb R$, $|p| \ge \mu(|z|)$, where $d_x = \dist(x,\partial\Omega)$,
$\mu$ is nondecreasing and $\beta = \gamma - n - 1 \ge 0$. Then we have the following gradient
estimate.

\textbf{Theorem 17.21.} \emph{Let $u \in C^0(\bar\Omega) \cap C^2(\Omega)$, $\phi \in C^2(\Omega) \cap C^0(\bar\Omega)$ be convex functions
in the uniformly convex domain $\Omega$, satisfying}

\begin{equation}
\tag{17.78}
\det D^2u = f(x,u,Du) \quad \text{in } \Omega, \qquad u = \phi \quad \text{on } \partial\Omega.
\end{equation}

\emph{Then we have}

\begin{equation}
\tag{17.79}
\sup_{\Omega} |Du| \le C,
\end{equation}

\emph{where $C$ depends on $n$, $\mu$, $\beta$, $\mathcal N$, $\Omega$, $|u|_{0;\Omega}$ and $|\phi|_{1;\Omega}$.}

\emph{Proof.} Let $B = B_R(y)$ be an enclosing sphere for the domain $\Omega$ at the point $x_0 \in \partial\Omega$
and set

\[
w = \phi - \psi(d)
\]

where $d(x) = \dist(x,\partial B)$ and $\psi$ is given by (14.11), with $v$ and $k$ to be determined.
Using a principal coordinate system for $\partial B$ at $x_0$, we may then estimate

\[
\det D^2w \ge \det(-D^2\psi)
\]

\[
= -\psi''\left(\frac{\psi'}{|x-y|}\right)^{n-1}
\]

\[
\ge -\psi''\left(\frac{\psi'}{R}\right)^{n-1}.
\]

% PDF page 485
\source{gilbarg-trudinger:page-0485}


while from the structure condition (17.77) we have

\[
f(x,u(x),Dw) \le 2^\gamma \mu(M)d^\beta(\psi')^\gamma
= 2^\gamma \mu(M)(\psi')^{n+1}(d\psi')^\beta
\]

provided $x \in \N$, $\psi'(d) \ge \mu(M)+|D\phi|$, where $M = |u|_{0;\Omega}$. Choosing $v =
1 + 2^\gamma R^{n-1}\mu$, so that $d\psi' \le 1$, and then choosing $k$ and $a$ according to (14.14) (with
$\mu$ replaced by $\mu + |D\phi|$), and such that $\{x \in \Omega| d < a\} \subset \N$, we see that the convex
function $w$ will be a lower barrier at $x_0$ for (17.63) and function $u$. Consequently, by
the comparison principle (Theorem 17.1), we obtain

\[
\frac{u(x) - u(x_0)}{|x - x_0|} \ge -C
\]

for $d \le a$, where $C$ depends on $n$, $\beta$, $\mu$, $|D\phi|_{0;\Omega}$, $M$, $\N$ and $R$. Using the convexity of
$u$, we have

\[
\frac{u(x) - u(x_0)}{|x - x_0|} \le |D\phi|_{0;\Omega},
\]

for all $x \in \Omega$, $x_0 \in \partial\Omega$, and the estimate (17.79) follows. $\square$

Note that the Monge-Ampère equation (17.2) is covered by Theorem 17.21 for
bounded $f$, while the equation of prescribed Gauss curvature is encompassed only
if the curvature $K$ vanishes Lipschitz continuously on $\partial\Omega$. Combining Theorems
17.4, 17.8, 17.10, 17.20 and 17.21, we obtain the following existence theorem for
equations of Monge-Ampère type in two variables.

\textbf{Theorem 17.22.} \emph{Let $\Omega$ be a uniformly convex domain in $\mathbb R^2$ with boundary $\partial\Omega \in C^3$.
Suppose that $f$ is a positive function in $C^2(\bar\Omega \times \mathbb R \times \mathbb R^n)$ satisfying $f_z \ge 0$, together
with the structure conditions (17.13), (17.14) and (17.77) for $n = 2$, $\beta = 0$. Then the
classical Dirichlet problem}

\begin{equation}
\tag{17.80}
\det D^2u = f(x,u,Du) \quad \text{in } \Omega, \qquad u = 0 \quad \text{on } \partial\Omega,
\end{equation}

\emph{is solvable with solution $u \in C^{2,\alpha}(\bar\Omega)$ for all $\alpha < 1$.}

\emph{Proof.} By virtue of Theorem 17.8 it suffices to have a uniform estimate in the space
$C^{2,\alpha}(\bar\Omega)$, for some $\alpha > 0$, of the solutions of the Dirichlet problems,

\[
F[u] = \frac{\det D^2u}{f(x,u,Du)} - 1 = \sigma F[\phi], \qquad u = 0 \quad \text{on } \partial\Omega,\quad 0 \le \sigma \le 1,
\]

where $\phi \in C^2(\bar\Omega)$ is uniformly convex and vanishes on $\partial\Omega$; (see Problem 17.8).
When the quantity $g_\infty$ in Theorem 17.4 is infinite, such an estimate follows im-
mediately from Theorems 17.4, 17.10, 17.20 and 17.21. Otherwise, we should, by
possibly a new choice of $\phi$, ensure that $F[\phi] \le 0$; (for example by replacing $\phi$ with
the solution of the Dirichlet problem, $\det D^2u = \inf f$, $u = 0 \quad \text{on } \partial\Omega$). $\square$

For Monge-Ampère equations in higher dimensions we have the following
analogue of Theorem 17.22, as a consequence of the global estimate (Theorem
17.26) in Section 17.8.

% PDF page 486
\source{gilbarg-trudinger:page-0486}


\textbf{Theorem 17.23.}\quad Let $\Omega$ be a uniformly convex domain in $\mathbb{R}^n$ with boundary $\partial\Omega\in C^4$.
Suppose that $f$ is a positive function in $C^2(\overline{\Omega}\times\mathbb{R}\times\mathbb{R}^n)$ satisfying $f_z\geq 0$ together
with the structure conditions (17.13), (17.14) and (17.77) for $\beta=0$. Then the classical
Dirichlet problem (17.80) is solvable with solution $u\in C^{3,\alpha}(\overline{\Omega})$ for all $\alpha<1$.

\medskip

\noindent We note here that Theorems 17.20, 17.22 and 17.23 extend to general boundary
values $\varphi\in C^4(\overline{\Omega})$; (see [IC 2], [CNS]).

\medskip

Using the interior second-derivative estimates of Theorem 17.20, the above
existence theorems can be extended to allow more general functions $f$.

\medskip

\textbf{Theorem 17.24.}\quad Let $\Omega$ be a uniformly convex domain in $\mathbb{R}^n$ and suppose that $f$ is a
positive function in $C^2(\Omega\times\mathbb{R}\times\mathbb{R}^n)$ satisfying $f_z\geq 0$ together with the structure
conditions (17.13), (17.14) and (17.77). Then the classical Dirichlet problem (17.80) is
solvable with solution $u\in C^{3,\alpha}(\Omega)\cap C^{0,1}(\overline{\Omega})$ for all $\alpha<1$.

\medskip

\noindent\textit{Proof.}\quad Let $\{f_m\}$ be a sequence of bounded functions in $C^2(\Omega\times\mathbb{R}\times\mathbb{R}^n)$ satisfying
$0<f_m\leq f$, $f_z\geq 0$ and $f_m=f$ for $|z|+|p|\leq m$ and let $\{\Omega_l\}$ be an increasing
sequence of uniformly convex $C^4$ subdomains of $\Omega$ satisfying $\Omega_l\Subset\Omega$, $\bigcup \Omega_l=\Omega$.
By Theorems 17.22 and 17.23, there exists, for each $m$, a sequence $\{u_{ml}\}$ of uniformly
convex solutions of the Dirichlet problems.

\[
\det D^2u_{ml}=f_m(x,u_{ml},Du_{ml})\quad\text{in }\Omega_l,\qquad u_{ml}=0\quad\text{on }\partial\Omega_l.
\]

Using the bounds, Theorems 17.4, 17.21 and the interior estimates, Theorems 17.14,
17.19, we obtain a subsequence converging uniformly, together with its first and
second derivatives on compact subsets of $\Omega$, to a solution $u_m$ of the Dirichlet problem

\[
\det D^2u_m=f_m(x,u_m,Du_m)\quad\text{in }\Omega,\qquad u_m=0\quad\text{on }\partial\Omega.
\]

But by Theorems 17.4, 17.21 again, it follows that for large enough $m$, $u_m=u$ is a
solution of (17.80). \hfill $\square$

\medskip

For the special case of the equation of prescribed Gauss curvature (17.3) we now
obtain from Theorem 17.24.

\medskip

\textbf{Corollary 17.25.}\quad Let $\Omega$ be a uniformly convex domain in $\mathbb{R}^n$ and $K$ a positive function
in $C^2(\Omega)\cap C^{0,1}(\overline{\Omega})$, satisfying

\[
\text{(17.81)}\qquad K=0\quad\text{on }\partial\Omega,\qquad \int_\Omega K<\omega_n.
\]


% PDF page 487
\source{gilbarg-trudinger:page-0487}

Then there exists a unique convex function $u \in C^2(\Omega) \cap C^{0,1}(\overline{\Omega})$ such that $u = 0$ on $\partial \Omega$ and whose graph has Gauss curvature $K(x)$ at each point $x \in \Omega$.

Both conditions (17.81) in Corollary 17.25 are necessary in some sense. In fact, let us suppose that in the general Monge-Ampère equation (17.4) the function $f$ satisfies

\begin{equation}
\tag{17.82}
f(x, z, p) \geq \frac{h(x)}{g(p)} \qquad \forall (x, z, p) \in \Omega \times \mathbb{R} \times \mathbb{R}^n,
\end{equation}

where $h$ and $g$ are positive functions in $L^1(\Omega)$, $L^1_{\mathrm{loc}}(\mathbb{R}^n)$, respectively. Then if $u \in C^2(\Omega)$ is a convex solution of (17.4) in a domain $\Omega$, its normal mapping $\chi$ coincides with $Du$ and is one-to-one. Consequently, by integration we have

\[
\int_{\chi(\Omega)} g(p)\, dp = \int_{\Omega} g(Du)\, \det D^2u
\geq \int_{\Omega} h,
\]

and the condition

\begin{equation}
\tag{17.83}
\int_{\Omega} h \leq g_{\infty} = \int_{\mathbb{R}^n} g(p)\, dp
\end{equation}

is thus necessary for the existence of a convex solution $u$. Furthermore the strict inequality

\begin{equation}
\tag{17.84}
\int_{\Omega} h < g_{\infty}
\end{equation}

is then necessary for the existence of a solution $u$ whose normal mapping is not all of $\mathbb{R}^n$, in particular a solution $u \in C^{0,1}(\overline{\Omega})$.

Concerning the other condition in (17.81) we remark that by extension of the interior estimate (Theorem 17.19), we can permit arbitrary nonzero boundary values $\phi \in C^2(\overline{\Omega})$ in Corollary 17.25 if and only if $K = 0$ on $\partial \Omega$ [TU]. The conditions (17.81) are analogous to the conditions (16.58), (16.60) for the equation of prescribed mean curvature (16.1).

% PDF page 488
\source{gilbarg-trudinger:page-0488}


17.8. Global Second Derivative Hölder Estimates

We consider in this section the global analogue, for equations of Monge-Ampère type, of the interior second-derivative Hölder estimates of Theorem 17.14. The methods automatically embrace equations of the general form (17.1) provided we strengthen hypothesis (ii)' in Theorem 17.14 by requiring

\begin{equation}
\tag{17.85}
-F_{ij,kl}(x,u,Du,D^2u)\xi_{ij}\xi_{kl} \geq \lambda_0|\xi|^2
\end{equation}

for some positive constant $\lambda_0$ and all symmetric matrices $\xi \in \R^{n \times n}$. By (17.65) we see that the equations of Monge-Ampère type, when written in the form (17.64), satisfy (17.85) with $\lambda_0 = (\sup \mathcal{C}_n)^{-2}$, where $\mathcal{C}_n$ denotes the maximum eigenvalue of $D^2u$. Since the hypotheses for the global second-derivative Hölder estimation guarantee, through the linear theory, stronger third-derivative estimates, it is convenient to formulate the results accordingly as follows.

\textbf{Theorem 17.26.} \emph{Let $\Omega$ be a bounded domain in $\R^n$ with boundary $\partial\Omega \in C^4$ and let $\phi \in C^4(\overline{\Omega})$. Suppose that $u \in C^3(\overline{\Omega}) \cap C^4(\Omega)$ satisfies $F[u] = 0$ in $\Omega$, $u = \phi$ on $\partial\Omega$, where $F \in C^2(\overline{\Gamma})$, $F$ is elliptic with respect to $u$ and satisfies (i)', (ii)' together with (17.85). Then we have the estimate,}

\begin{equation}
\tag{17.86}
\sup_{\Omega} |D^3u| \leq C,
\end{equation}

\emph{where $C$ depends on $n$, $\lambda$, $\Lambda$, $\lambda_0$, $\partial\Omega$, $|\phi|_{4;\Omega}$, $|u|_{2;\Omega}$ and the first and second derivatives of $F$.}

\emph{Proof.} Letting $\gamma$ denote a unit vector in $\R^n$, we can, by use of (17.85) in (17.43), improve the inequality (17.44) so that

\begin{equation}
\tag{17.87}
F_{ij}D_{ij\gamma\gamma}u \geq -A_{ij\gamma}D_{ij\gamma}u - B_{\gamma} + \lambda_0|D^2D_{\gamma}u|^2
\geq -C
\end{equation}

by Cauchy's inequality, where $C$ depends on $n$, $\lambda_0$ and $A_0$, $B_0$ in (17.46). Let us now fix a point $x_0 \in \partial\Omega$ and suppose that $B_R(y)$ is an exterior ball at $x_0$. The modulus of continuity of the derivatives $D_{yy}u$ can be estimated at the boundary (in terms of their traces on $\partial\Omega$) by standard barrier arguments, as described for example in Remark 3 in Section 6.3. As shown there, the function $w$ in (6.45) given by

\[
w(x)=\tau(R^{-\sigma}-|x-y|^{-\sigma})
\]

satisfies

\[
F_{ij}D_{ij}w \leq -1
\]

in $\Omega$ for sufficiently large $\tau$, $\sigma$ depending on $\lambda$, $\Lambda$ and $R$. Consequently, if $\epsilon \geq 0$ and

\[
|D_{yy'}u(x)-D_{yy'}u(x_0)| \leq \epsilon
\]

% PDF page 489
\source{gilbarg-trudinger:page-0489}

for $x \in \partial \Omega$, $|x - x_0| < \delta$, we obtain from (17.87) and the classical maximum principle (Theorem 3.1), the estimate

\begin{equation}
\tag{17.88}
D_{\gamma \gamma} u(x) - D_{\gamma \gamma} u(x_0) \leq \varepsilon + C w + 2 \sup_{\partial \Omega} |D_{\gamma \gamma} u| \, |x - x_0|^2 / \delta^2
\leq \varepsilon + C |x - x_0|
\end{equation}

for any $x \in \Omega$, where $C$ depends on $n$, $\lambda$, $\Lambda$, $\lambda_0$, $A_0$, $B_0$, $\|u\|_{2;\Omega}$, diam $\Omega$, $R$ and $\delta$. Now choosing directions $\gamma_1, \ldots, \gamma_N$ in accordance with Lemma 17.13 and using equation (17.1) itself, as in the proof of Theorem 17.14, we conclude from (17.88), the estimate

\begin{equation}
\tag{17.89}
\bigl|D^2 u(x) - D^2 u(x_0)\bigr| \leq \varepsilon + C |x - x_0|
\end{equation}

for any $x \in \Omega$, where $C$ depends in addition on the quantity $D_0$ in the proof of Theorem 17.14.

The estimate (17.89) reduces estimation of the modulus of continuity of second derivatives on the boundary $\partial \Omega$ to that of their traces on $\partial \Omega$. A similar result could also have been established by using the weak Harnack inequality at the boundary (Theorem 9.27), and moreover can be proved without requiring the condition (17.85).

To proceed further we use (17.88) again to obtain a one-sided third derivative bound. Without loss of generality we may assume that $u$ vanishes on $\partial \Omega$, and that $\partial \Omega$ is flat in a neighbourhood of $x_0 \in \partial \Omega$, that is, for some $\delta > 0$,

\[
B^+ = \Omega \cap B_\delta(x_0) \subset \mathbb{R}^n_+,
\]

\[
T = \partial \Omega \cap B_\delta(x_0) \subset \partial \mathbb{R}^n_+.
\]

This follows since the form of (17.1), together with (17.85) and the hypotheses (i), (ii) in Theorem 17.8, are preserved by replacement of $u$ with $u - \phi$ and by a $C^4$ coordinate change $\psi$. The new constants $\lambda$, $\Lambda$ and $\lambda_0$ will, of course, depend on $\psi$ (in particular on $D\psi$) as well as their original values. Therefore, restricting $\gamma$ to tangential directions in $\partial \mathbb{R}^n_+$ so that $D_{\gamma \gamma} u = 0$ on $T$, we have from (17.88)

\[
D_{\gamma \gamma} u(x) - D_{\gamma \gamma} u(x_0) \leq C |x - x_0|
\]

for $x \in B^+$; and consequently

\begin{equation}
\tag{17.90}
D_{\gamma \gamma n} u(x_0) \geq -C,
\end{equation}

where $C$ depends on $\|\phi\|_{4;\Omega}$ and $\psi$, as well as the quantities in (17.88). Setting

\[
h = D_n u + k |x|^2
\]

it follows that for sufficiently large $k$ ($k \geq C$),

\[
D_{\gamma \gamma} h(x) \geq 0
\]

% PDF page 490
\source{gilbarg-trudinger:page-0490}

for $x \in T$, $|x - x_0| \leq \delta/2$; that is, the function $h$ has convex trace on $\partial \Omega \cap B_{\delta/2}(x_0)$. Furthermore, we see from the differentiated equation (17.23) that $h$ satisfies a uniformly elliptic equation, in particular

\[
|F_{ij} D_{ij} h| \leq C
\]

where $C$ depends on the quantities in the statement of Theorem 17.26. The above properties of the function $h$ are utilized through the following remarkable lemma, whose proof we defer until the end of this section.

\textbf{Lemma 17.27.} \emph{Let $h \in C^2(B^+) \cap C^0(B^+ \cup T)$ satisfy}

\begin{equation}
\tag{17.91}
Lh = a^{ij}D_{ij}h \leq f
\end{equation}

\emph{in $B^+$ where $L$ is uniformly elliptic in $B^+$ and $f/\lambda$ is bounded. Then if $h|_T$ is convex, we have for any $x$, $y \in T \cap B_{\delta/2}(x_0)$ and $i=1,\ldots,n-1$, the estimate}

\begin{equation}
\tag{17.92}
|D_i h(x) - D_i h(y)| \leq \frac{C}{1 + |\log |x-y||}
\end{equation}

\emph{where $C$ depends only on $n$, sup $\Lambda/\lambda$, sup $f/\lambda$, $\delta$ and sup $|Dh|$.}

Lemma 17.27 provides an estimate for the moduli of continuity at $x_0$ of the second derivatives $D_{in}u$ restricted to $T$. A similar estimate for the remaining non-tangential second derivative $D_{nn}u$ follows immediately from (17.1) itself. But then the estimate (17.89) yields an estimate for the modulus of continuity of the full Hessian matrix $D^2u$ at the boundary point $x_0$ in terms of the quantities specified in the statement of Theorem 17.26, which in turn implies an estimate for the moduli of continuity of the principal coefficients of the differentiated equation (17.23). We can then use the $L^p$ theory, in particular Theorem 9.13, to conclude $L^p$ estimates for the third derivatives $D_{ijk}u$, $i,j=1,\ldots,n$, $k=1,\ldots,n-1$, in a neighbourhood of the point $x_0$ for any $p < \infty$. Similar estimates follow then for the remaining third derivative $D_{nnn}u$ from (17.23) again. Using Morrey's estimate (Theorem 7.19), we thus obtain Hölder estimates for $D^2u$ at $x_0$ for any exponent $\alpha < 1$. On combination with the interior estimate (Theorem 17.14), as in the proof of Theorem 8.29, we finally infer global estimates in $C^{2,\alpha}(\overline{\Omega})$ for any $\alpha < 1$. To complete the proof of Theorem 17.26 we simply apply the Schauder theory to the differentiated equation (17.23), thereby obtaining global $C^{3,\alpha}(\overline{\Omega})$ estimates for any $\alpha < 1$. $\square$

\emph{Proof of Lemma 17.27.} Since $h|_T$ is convex the derivatives $D_i h$, $i=1,\ldots,n-1$, exist almost everywhere on $T$ and it suffices to prove (17.92) for such points $x$, $y$. The full result for all $x$, $y$ follows from convexity and continuity.

We may assume $x=0$ and that after subtraction of a suitable affine function,

\begin{equation}
\tag{17.93}
h(0)=0 = D_i h(0), \qquad i=1,\ldots,n-1.
\end{equation}

% PDF page 491
\source{gilbarg-trudinger:page-0491}

After a rotation of coordinates we may further assume that

\[
D_1 h(y) = \alpha \geq 0, \qquad D_i h(y) = 0, \qquad i = 2,\ldots,n-1,
\]

where $\alpha \leq \sup |Dh| \leq C$. In the following we shall use the same letter $C$ to denote constants depending only on the constants in the statement of the lemma. We wish to prove

\begin{equation}
\tag{17.94}
\alpha \leq \frac{C}{|\log |y||}
\end{equation}

when $|y|$ is sufficiently small, from which (17.92) will follow. Since $\alpha = 0$ implies the desired conclusion we may suppose $\alpha > 0$.

Writing $x = (x', x_n) = (x_1,\ldots,x_{n-1},x_n)$ for $x \in \mathbb{R}^n$, we have from the convexity of $h$ on $T$ that $h(x',0) \geq 0$ and

\begin{equation}
\tag{17.95}
h(x',0) \geq h(y',0) + \alpha(x_1 - y_1) = \alpha(x_1 - \beta),
\end{equation}

where $\beta$ is defined by the equality. By taking $x' = 0$ and $x' = y'$ in this relation, one sees that

\begin{equation}
\tag{17.96}
0 \leq \beta \leq y_1.
\end{equation}

We consider the function

\begin{equation}
\tag{17.97}
w(x',x_n) = \frac{\alpha}{2} \left[(x_1 - \beta)^2 + x_n^2\right]^{1/2} + \frac{\alpha}{2}(x_1 - \beta)
\end{equation}
\[
\qquad\qquad - \alpha \gamma x_n \log \left[(x_1 - \beta)^2 + x_n^2\right]^{1/2} - D(x_n + |x'|^2 - Ex_n^2).
\]

After suitable choice of positive constants $\gamma$, $\varepsilon < 1$, and $D$, $E > 1$, depending only on the constants of the lemma, it will be seen that $w$ provides a barrier with the properties

\begin{equation}
\tag{17.98}
Lw \geq f \geq Lh \qquad \text{in } B_\varepsilon^+(= B_\varepsilon^+(0))
\end{equation}

and

\begin{equation}
\tag{17.99}
w \leq h \qquad \text{on } \partial B_\varepsilon^+.
\end{equation}

Assuming for the moment that the constants in $w$ can be so chosen, we infer from the maximum principle that

\[
w \leq h \qquad \text{in } B_\varepsilon^+.
\]

% PDF page 492
\source{gilbarg-trudinger:page-0492}

Then since $h(0) = w(0) = 0$, we have, after setting $x' = 0$ in (17.97), dividing by $x_n$
and letting $x_n \to 0$,

\[
D_n w(0) \leq \liminf_{x_n \to 0} \frac{h(0,x_n)}{x_n} \leq \sup |Dh| \leq C,
\]

or

\[
-\alpha \gamma \log \beta - D \leq C.
\]

It follows from (17.96) that

\[
\alpha \leq \frac{C + D}{\gamma |\log \beta|} \leq \frac{C + D}{\gamma |\log |y'||}
\]

proving (17.94).
It remains to determine the constants in $w$. Setting $\rho = \left[(x_1 - \beta)^2 + x_n^2\right]^{1/2}$, we
obtain by direct calculation

\[
L \rho \geq \frac{\lambda}{\rho}
\]

and

\[
|L(x_n \log \rho)| \leq \frac{6\Lambda}{\rho},
\]

so that for $0 < \gamma \leq \frac{1}{12} \inf \lambda/\Lambda$ we have

\[
L(\rho/2 - \gamma x_n \log \rho) \geq 0.
\]

By now choosing the constant $E$ sufficiently large $(\geq \frac{1}{2} \sup f/\lambda + n \sup A/\lambda)$ and
$D \geq 1$, we can satisfy (17.98).
To establish (17.99) we observe first that the inequality holds on $x_n = 0$. For

\[
w(x',0) = \alpha(x_1 - \beta)^+ - D|x'|^2
\]
\[
= \alpha(x_1 - \beta) - D|x'|^2 \leq h(x',0) \qquad \text{if } x_1 \geq \beta \qquad \text{(by (17.95))}
\]
\[
= -D|x'|^2 \leq 0 \leq h(x',0) \qquad \text{if } x_1 \leq \beta.
\]

On the hemispherical portion $S$ of $\partial B_\varepsilon^+$ we have that $x_n + |x'|^2 - E x_n^2 \geq \frac{1}{2}\varepsilon^2$ pro-
vided $\varepsilon$ is sufficiently small (depending on the choice of $E$). Since the other terms in
$w$ are bounded by $\sup |Dh|$ if $\varepsilon < 1$, a suitably large value of $D$ makes

\[
w \leq -C \leq \inf h
\]

on $S$, which completes the proof of (17.94).

% PDF page 493
\source{gilbarg-trudinger:page-0493}

Finally, we remark that \((17.94)\) implies \((17.92)\) if \(|x-y| \le \varepsilon\), while \(|Dh| \le C\)
yields the same inequality if \(|x-y| > \varepsilon\) in \(B_{\delta/2}^+\). \(\square\)

The proof of the above lemma is taken almost directly from [CNS].

\section*{An Alternative Approach}

The approach to global second derivative Hölder estimates given above is due to
Caffarelli, Nirenberg and Spruck [CNS]. However, as we shall indicate now, the
crucial Hölder estimates for the mixed tangential normal second derivatives on
the boundary also follow readily from Theorem 9.3.1. Furthermore this approach,
due to Krylov [KV 5], also yields global regularity for the Bellman equation,
(17.8). In fact we have the stronger version of Theorem 17.6.

\textbf{Theorem 17.26\(^\prime\).} \textit{Let \(\Omega\) be a bounded domain in \(\mathbb{R}^n\) with boundary \(\partial\Omega \in C^3\) and let
\(\phi \in C^3(\overline{\Omega})\). Suppose that \(u \in C^3(\overline{\Omega}) \cap C^4(\Omega)\) satisfies \(F[u] = 0\) in \(\Omega\), \(u = \phi\) on \(\partial\Omega\),
where \(F \in C^2(\Gamma)\), \(F\) is elliptic with respect to \(u\) and satisfies (i)', (ii)'. Then we have
the estimate}

\[
(17.86)\qquad |u|_{2,\alpha;\Omega} \le C
\]

\textit{where \(\alpha\) depends only on \(n\), \(\lambda\) and \(\Lambda\) and \(C\) depends in addition on \(|u|_{2;\Omega}\) and the first
and second derivatives of \(F\) other than \(F_{rr}\).}

\textit{Proof.} The situation is reduced according to the proof of Theorem 17.26 to the
consideration of the derivatives \(D_i u\), \(i = 1,\ldots,n-1\) on the flat boundary por-
tion \(T\). But now, instead of considering the differential inequality \((17.91)\) satisfied
by the normal derivative \(D_n u\) and applying Lemma 17.27, we consider the uni-
formly elliptic differential equations

\[
a^{ij} D_{ij} h = f_k
\]

satisfied by the tangential derivatives \(h = D_k u, k = 1,\ldots,n-1\) in \(B^+\) and apply
Theorem 9.31. A global Hölder estimate for \(D^2 u\) then follows independently of the
moduli of continuity of the principal coefficients of the differentiated equation
(17.23). \(\square\)

We remark here, that by using the particular forms of the estimates (9.68),
(17.51), we can avoid the barrier argument at the beginning of the proof of Theo-
rem 17.26; (see Problem 13.1 with \(u\) replaced by \((D_1 u, \ldots, D_{n-1} u)\) or [TR 14]).
Also the present proof of Theorem 17.26 can be modified slightly so as to remove
the restriction (17.85). To see this, we recall that the pure second derivatives,
\(v = D_{yy} u = \gamma_i \gamma_j D_{ij} u\), \(\gamma \in \partial \mathbb{R}_+^n\), satisfy in \(B^+\), uniformly elliptic differential in-
equalities of the form

\[
F_{ij} D_{ij} v + C_{ij} D_{iyy} u \ge -C
\]

where by virtue of the differentiated equation (17.23), we may assume \(C_{in} = 0\),
\(i = 1,\ldots,n\). The coefficients \(C_{ij}\) and constant \(C\) will be bounded in terms of the

% PDF page 494
\source{gilbarg-trudinger:page-0494}

quantities in the statement of Theorem 17.26 (excluding $\lambda_0$). Now let us regard $v$ as
a function of both $x$ and $\gamma$ in the domain $\Omega^* = B^+ \times \{ \gamma \in \partial \mathbb{R}_+^n \mid |\gamma| < 2\} \subset \mathbb{R}^{2n-1}$
and extend the above operator so that

\[
\tilde{L}v \equiv F_{ij} D_{ij} v + \frac{1}{2} C_{ij} D_i v + K_0 \sum_{i=1}^{n-1} D_{\gamma_i \gamma_i} v \geq -C
\]

in $\Omega^*$ and $K_0$ is chosen so that $\tilde{L}$ is uniformly elliptic in $\Omega^*$. By appropriate
extension of the barrier $w$ to $\Omega^*$, for example by defining

\[
w(x,\gamma) = w(x) + \tau' |\gamma - \bar{\gamma}|^2
\]

for constant $\tau'$ and $|\bar{\gamma}| = 1$, we are again able to infer the one-sided third derivative
bound (17.90), and the rest of the proof of Theorem 17.26 follows automatically.
    As a consequence of Theorem 17.26 we can deduce global smoothness of the
solutions of the Dirichlet problems in Theorems 17.17, 17.18 when the boundary
data are appropriately smooth. In fact, if $F$ satisfies the structure conditions
(17.53), we can, by interpolation, replace the dependence of the constant $C$ in
(17.86)' on $\|u\|_{2;\Omega}$ by $\|u\|_{0;\Omega}$, and furthermore this estimate will be uniform with
respect to the approximation of the Dirichlet problem (17.61) treated in Sec-
tion 17.5. We thus obtain, for example, in place of Theorem 17.18.

\textbf{Theorem 17.18'.} \textit{Let $\Omega$ be a bounded domain in $\mathbb{R}^n$ with boundary $\partial\Omega \in C^3$ and let
$\phi \in C^3(\overline{\Omega})$. Suppose the functions $F^1, F^2, \ldots \in C^2(\overline{\Gamma})$, are concave with respect to
$z, p, r$, non-increasing with respect to $z$, and satisfy uniformly the structure conditions
(17.53). Then the classical Dirichlet problem, (17.61), is uniquely solvable with solu-
tion $u \in C^{2,\alpha}(\overline{\Omega})$ for some positive $\alpha$ depending only on $n$, $\lambda$ and $\Lambda$.}

\section*{17.9. Nonlinear Boundary Value Problems}

The method of continuity, previously applied to the Dirichlet problem in Section
17.2, readily extends to other boundary value problems as well. Even for the quasi-
linear case, the fixed point methods of Chapter 11 are not appropriate, since it is not
in general possible to construct a \textit{compact} operator analogous to the operator $T$ in
Sections 11.2, 11.4.
    We first consider nonlinear boundary value problems of the form,

\[
\begin{aligned}
F[u] &= F(x,u,Du,D^2u)=0 && \text{in } \Omega, \\
G[u] &= G(x,u,Du) && = 0 && \text{on } \partial\Omega,
\end{aligned}
\tag{17.100}
\]

where $F$ and $G$ are real functions on the sets $\Gamma = \Omega \times \mathbb{R} \times \mathbb{R}^n \times \mathbb{R}^{n\times n}$ and
$\Gamma' = \partial\Omega \times \mathbb{R} \times \mathbb{R}^n$ respectively. The case

\[
G(x,z,p)= z - \varphi(x)
\tag{17.101}
\]

corresponds to the Dirichlet problem, $F[u] = 0$ in $\Omega$, $u = \varphi$ on $\partial\Omega$, previously
treated in this work. If $\partial\Omega \in C^1$, $G \in C^1(\Gamma')$, the boundary operator $G$ is called

% PDF page 495
\source{gilbarg-trudinger:page-0495}

oblique if

\begin{equation}
\tag{17.102}
\frac{\partial G}{\partial p}(x,z,p)\cdot v(x) > 0
\end{equation}

for all $(x,z,p)\in \Gamma'$, where $v$ denotes the outer unit normal to $\partial\Omega$, while if

\begin{equation}
\tag{17.103}
\frac{\partial G}{\partial p}(x,u(x),Du(x))\cdot v(x) > 0
\end{equation}

for all $x\in \partial\Omega$, for some function $u\in C^{1}(\overline{\Omega})$ we call $G$ oblique with respect to $u$. To
apply the method of continuity we shall assume that $F\in C^{2,\alpha}(\overline{\Gamma})$, $G\in C^{3,\alpha}(\Gamma')$ and
take for our Banach spaces,

\[
\mathbb{B}_1 = C^{2,\alpha}(\overline{\Omega}), \mathbb{B}_2 = C^{\alpha}(\overline{\Omega}) \times C^{1,\alpha}(\partial\Omega)
\]

for some $\alpha \in (0,1)$. We then define a mapping $P : \mathbb{B}_1 \to \mathbb{B}_2$ by

\[
P[u] = (F[u], G[u]),
\]

which has a continuous Fr\'echet derivative on $\mathbb{B}_1$ given by

\begin{equation}
\tag{17.104}
P_u h = (Lh, Nh), \qquad h\in \mathbb{B}_1,
\end{equation}

where $L$ is defined by (17.22) and

\[
Nh = \gamma(x)h + \beta^i(x)D_i h
\]

with

\[
\gamma(x) = G_z(x,u(x),Du(x)),
\]
\[
\beta^i(x) = G_{p_i}(x,u(x),Du(x)).
\]

It follows then, from the Schauder theory for the linear oblique derivative problem
in Theorem 6.31, that $P_u$ is boundedly invertible if $F$ is strictly elliptic with respect
to $u$, $G$ is oblique with respect to $u$, $c(=L1)\leq 0$, $\gamma \geq 0$ and either $c \neq 0$ in $\Omega$ or
$\gamma \neq 0$ on $\partial\Omega$. Accordingly we have, by virtue of Theorem 17.7, the following exten-
sion of Theorem 17.8.

\textbf{Theorem 17.28.} Let $\Omega$ be a $C^{2,\alpha}$ domain in $\mathbb{R}^n$, $0 < \alpha < 1$, $\mathcal{U}$ an open subset of the
space $C^{2,\alpha}(\overline{\Omega})$ and $\psi \in \mathcal{U}$. Set

\begin{equation}
\tag{17.105}
E = \{u \in \mathcal{U}\mid P[u] = \sigma P[\psi] \text{ for some } \sigma \in [0,1]\}.
\end{equation}

Suppose that $F \in C^{2,\alpha}(\overline{\Gamma})$, $G \in C^{3,\alpha}(\Gamma')$ and that

\begin{enumerate}
\item[(i)] $F$ is strictly elliptic in $\Omega$ for each $u \in E$;
\item[(ii)] $G$ is either of the form (17.101) or oblique on $\partial\Omega$ for each $u \in E$;
\end{enumerate}


% PDF page 496
\source{gilbarg-trudinger:page-0496}

(iii) $F_z(x, u, Du, D^2u) \leq 0$, $G_z(x, u, Du) \geq 0$ for each $u \in E$ with one of these
quantities not vanishing identically;
(iv) $E$ is bounded in $C^{2,\alpha}(\overline{\Omega})$;
(v) $\overline{E} \subset \mathcal{U}$.

Then the boundary value problem (17.100) is solvable in $\mathcal{U}$.

    A remark analogous to that at the end of Section 17.2 is also pertinent here.
Namely, the family of boundary value problems in (17.105) may be replaced by
more general families, $P[u;\sigma] = 0$, depending smoothly on $\sigma$, for which $P[u; 1] =
P[u]$ and the equation $P[u; 0] = 0$ is solvable in $\mathcal{U}$. For each $\sigma \in [0, 1]$ the operator
$u \to P[u;\sigma]$ must of course satisfy the same hypotheses as $P$.
    We now show that for quasilinear operators

\[
(17.106)\qquad F[u] = Qu = a^{ij}(x, u, Du)D_{ij}u + b(x, u, Du),
\]

condition (iv) in Theorem 17.28 can be weakened to require only the boundedness
of the set $E$ in $C^{1,\delta}(\overline{\Omega})$ for some $\delta > 0$. This has the effect of reducing the existence
program for quasilinear equations to the same type of a priori estimation as was
required for the Dirichlet problem in Chapter 11. For this reduction we need the
following convergence lemma.

\textbf{Lemma 17.29.} \textit{Let $F$ be a quasilinear operator of the form (17.106) with coefficients
$a^{ij}$, $b \in C^1(\Omega \times \mathbb{R} \times \mathbb{R}^n)$ and $G$ a boundary operator of the form (17.100) with $G$,
$G_p \in C^{1,\alpha}(\partial\Omega \times \mathbb{R} \times \mathbb{R}^n)$, $0 < \alpha < 1$. Assume that $F$ is strictly elliptic in $\Omega$, $G$ is
oblique on $\partial\Omega$, and suppose there exist sequences $\{u_m\} \subset C^{2,\alpha}(\overline{\Omega})$, $\{f_m\} \subset C^\alpha(\overline{\Omega})$,
$\{\varphi_m\} \subset C^{1,\alpha}(\partial\Omega)$ such that}

\[
\begin{aligned}
&\text{(i)}\quad P[u_m] = (f_m, \varphi_m);\\
&\text{(ii)}\quad \{u_m\}, \{f_m\}, \{\varphi_m\} \text{ are uniformly bounded in } C^{1,\alpha}(\overline{\Omega}),\ C^\alpha(\overline{\Omega}),\ C^{1,\alpha}(\partial\Omega),\\
&\qquad\qquad\text{respectively;}\\
&\text{(iii)}\quad \{u_m\}, \{f_m\}, \{\varphi_m\} \text{ converge uniformly to functions } u, f, \varphi, \text{ respectively.}
\end{aligned}
\]

Then $u \in C^{2,\alpha}(\overline{\Omega})$ and $P[u] = (f, \varphi)$.

\textit{Proof.} \quad Let $v = u_m - u_k$ where $m, k$ are positive integers. Then $v$ is a solution of the
linear problem

\[
Lv = g \quad \text{in } \Omega, \qquad Nv = \psi \quad \text{on } \partial\Omega
\]

where

\[
Lu = \bar{a}^{ij}D_{ij}u,\qquad Nu = \beta^i D_i u,
\]

\[
\bar{a}^{ij} = a^{ij}(x, u_m, Du_m), \qquad \beta^i = \int_0^1 G_{p_i}(x, u_m, Du_k + t(Du_m - Du_k))\,dt,
\]

\[
g = f_m - f_k + b(x, u_k, Du_k) - b(x, u_m, Du_m)
\]
\[
\qquad + \bigl(a^{ij}(x, u_k, Du_k) - a^{ij}(x, u_m, Du_m)\bigr)D_{ij}u_k,
\]

\[
\psi = \varphi_m - \varphi_k + G(x, u_k, Du_k) - G(x, u_m, Du_m).
\]

% PDF page 497
\source{gilbarg-trudinger:page-0497}

Notes

From the hypotheses of the lemma we estimate

\[
|g|_{\alpha\delta} \le \epsilon(m,k)(1 + K_1 + K_2) + CK_1,
\]
\[
|\psi|_{1,\alpha\delta} \le \epsilon(m,k)(1 + K_1 + K_2)
\]
\[
|\beta|_{\alpha\delta} \le C,\ [\beta]_{1,\alpha\delta} \le C(1 + K_1 + K_2)
\]

where $\epsilon \to 0$ as $m,k \to \infty$, $C$ is independent of $m,k$ and

\[
K_1 = \max \{ |D^2u_m|_{0}, |D^2u_k|_{0} \}, \qquad
K_2 = \max \{ [D^2u_m]_{\alpha\delta}, [D^2u_k]_{\alpha\delta} \}.
\]

Using the estimate of Theorem 6.30, as modified in Problem 6.11, we thus obtain

\[
[v]_{2,\alpha\delta} \le C(1 + K_1) + \epsilon(m,k)K_2
\]

and hence by the interpolation inequality, Theorem 6.35, we have

\[
(17.107)\qquad [v]_{2,\alpha\delta} \le C + \left(\frac14 + \epsilon\right)K_2,
\]

where again $C$ is independent of $m,k$ and $\epsilon \to 0$ as $m,k \to \infty$. But (17.107) implies
the boundedness of $[u_m]_{2,\alpha\delta}$. For suppose not. Then for some $m,k$

\[
(17.108)\qquad K_2 = [u_m]_{2,\alpha\delta} > 4[u_k]_{2,\alpha\delta} \ge 4C,\qquad \epsilon(m,k) < \frac14.
\]

Hence

\[
[v]_{2,\alpha\delta} \ge \frac34 K_2,
\]

so that by (17.107) we have

\[
\frac34 K_2 \le C + \frac12 K_2,
\]

that is, $K_2 \le 4C$, contradicting (17.108). Thus $[u_m]_{2,\alpha\delta}$ is bounded independently of
$m$ and hence $u \in C^{2,\alpha\delta}(\overline{\Omega})$ with $\Pop[u] = (f,\varphi)$. Finally, by an argument similar to
the above we infer also $u \in C^{2,\alpha}(\overline{\Omega})$. $\square$

Combining Lemma 17.29 with Theorem 17.7 we now conclude a version of
Theorem 17.28 appropriate for quasilinear operators.

\textbf{Theorem 17.30.} \ Let $\Omega$ be a $C^{2,\alpha}$ domain in $\R^n$, $0 < \alpha < 1$, and suppose that $Q$ is
strictly elliptic in $\Omega$ with $a^{ij}$, $b \in C^2(\overline{\Omega}\times \R \times \R^n)$, and $G \in C^3(\partial\Omega \times \R \times \R^n)$ is
oblique on $\partial\Omega$. Suppose also that $a^{ij}_z=0$, $b_z \le 0$, $G_z \ge 0$ with either $b_z(x,u(x),Du(x))$
$\neq 0$ or $G_z(x,u(x),Du(x)) \neq 0$ for each $u \in C^{2,\alpha}(\overline{\Omega})$. Then if for some function $\psi \in
C^{2,\alpha}(\overline{\Omega})$, the set

\[
(17.109)\qquad E = \{u \in C^{2,\alpha}(\overline{\Omega}) \mid \Qop u = \sigma \Qop\psi \quad \text{in } \Omega,
\]
\[
G[u] = \sigma G[\psi] \quad \text{on } \partial\Omega \quad \text{for some } \sigma \in [0,1]\}
\]

% PDF page 498
\source{gilbarg-trudinger:page-0498}

is bounded in $C^{1,\delta}(\overline{\Omega})$ for some $\delta > 0$, the boundary value problem $Qu = 0$ in $\Omega$, $G[u] = 0$ on $\partial \Omega$ is uniquely solvable in $C^{2,\alpha}(\overline{\Omega})$.

As mentioned previously other families of problems may be used in (17.109). For example, we could consider a family analogous to that used for the Dirichlet problem in Theorem 11.4, namely

\begin{equation}
Q_{\sigma}u = a^{ij}(x,u,Du)D_{i j}u + \sigma b(x,u,Du) = 0 \qquad \text{in } \Omega,
\tag{17.110}
\end{equation}
\begin{equation}
G_{\sigma}u = \sigma G(x,u,Du) + (1-\sigma)u = 0 \qquad \text{on } \partial \Omega.
\end{equation}

A typical oblique boundary value problem is the conormal derivative boundary value problem for an elliptic divergence structure equation, which takes the form

\begin{equation}
Qu = \operatorname{div} A(x,u,Du) + B(x,u,Du) = 0 \qquad \text{in } \Omega
\tag{17.111}
\end{equation}
\begin{equation}
G[u] = A(x,u,Du)\cdot \nu(x) + \varphi(x,u) = 0 \qquad \text{on } \partial \Omega.
\end{equation}

The problem of determining a capillary surface with prescribed contact angle, mentioned in Chapter 10, provides an interesting special case of a conormal derivative boundary value problem. For this example and also for uniformly elliptic $Q$ satisfying natural conditions as in Theorem 15.11, the relevant apriori estimates have been proved, so that appropriate existence theorems may be inferred from Theorem 17.30. The reader is referred to the literature, for example [LU 4], [UR], [GE 3], [LB 2, 3] for further details.

\bigskip

\textbf{Notes}

Fully nonlinear equations in two variables were treated by various authors, including Lewy [LW 1], Nirenberg [NI 1], Pogorelov [PG 1] and Heinz [HE 2], with most attention being devoted to equations of Monge-Ampère type and related geometric problems such as the Minkowski problem to determine a convex hypersurface through prescription of its Gauss curvature. The extremal operators of Pucci were introduced in [PU 2] and the associated Dirichlet problems also solved in the two variable case.

The higher dimensional Monge-Ampère equations were first solved in a generalized sense by Aleksandrov [AL 1] and Bakelman [BA 1] using polyhedral approximation, with further development in [BA 2, 4], which contains a generalized version of the main existence theorem (Theorem 17.24). A generalized solution of the Monge-Ampère equation (17.2) can be defined as a convex function whose normal mapping is absolutely continuous with density $f$. The interior regularity of generalized solutions (under sufficiently smooth boundary conditions) was established by Pogorelov [PG 2, 3, 4, 5] and Cheng and Yau [CY 1, 2], essential ingredients of their proof being the interior second-derivative bounds of Pogorelov and the interior third-derivative bounds of Calabi [CL]. These authors thus obtained Theorem 17.24 for the Monge-Ampère equation (17.2). Our treatment of second-derivative estimates in Section 17.6 uses the Pogorelov method [PG 1, 5] (but incorporates a suggestion of L. M. Simon concerning the function $h$).


% PDF page 499
\source{gilbarg-trudinger:page-0499}

The main impetus for studying fully nonlinear uniformly elliptic equations
arose through stochastic control theory, as the Bellman equation (17.8) is satisfied
by a (sufficiently smooth) optimal cost function in a control problem associated
with a system of stochastic differential equations. The first significant treatment of
this equation by Krylov used probability methods and is described in his book
[KV 2]. Partial differential equation techniques were subsequently developed by:
Brezis and Evans (who derived $C^{2,\alpha}$ estimates for the two operator case); Evans and
Friedman [EF] for the constant coefficient case (see also [LD]); and P. L. Lions
[LP 4], Evans and Lions [EL 1] for the general uniformly elliptic case. In the
papers [LP 4], [EL 1], the existence of a strong solution of the Dirichlet problem is
established under conditions (17.62) using a clever method based on approxima-
tion by an elliptic system and apriori $C^{1,1}(\overline{\Omega})$ bounds. A full treatment of the
(possibly degenerate) Bellman equation and connections with stochastic control
theory and the Bellman dynamic programming method is given by P.-L. Lions in
[LP 8, 9] and discussed in [LP 10]; (see also [LP 7] for the first-order Hamilton-
Jacobi equation).

The theory of classical solutions of fully nonlinear elliptic equations advanced
substantially with the discovery of interior second-derivative H\"older estimates by
Evans [EV 2, 3] and Krylov [KV 4] who basically proved the results of Sec-
tions 17.4 and 17.5. These are treated in Sections 17.4, 17.5 following simplifica-
tions of Trudinger [TR 13]. While we have for expedience in Sections 17.3 and 17.4
deduced first- and second-derivative bounds through interpolation, such results
hold under more general structure conditions analogous to those of the quasilin-
ear theory [TR 13, 14]. In particular, we may replace the derivative conditions in
(17.29) and (17.53) by

\begin{align*}
(17.29)'\qquad &|F(x, z, p, 0)| \le \mu(1 + |p|^2),\\
&|F_x|, |F_z|, |F_p| \le \tilde{\mu}(1 + |r|);\\[0.5em]
(17.53)'\qquad &(1 + |p|)|F_{pp}|, |F_z|, |F_x| \le \mu(1 + |p|^2 + |r|),\\
&|F_{rx}|, |F_{rz}|, |F_{rp}| \le \tilde{\mu},\\
&|F_{xx}|, |F_{xz}|, |F_{xp}|, |F_{zz}|, |F_{zp}|, |F_{pp}| \le \tilde{\mu}(1 + |r|)
\end{align*}

where $\mu$ is nondecreasing in $|z|$, $\tilde{\mu}$ is nondecreasing in $|z| + |p|$, (and the concavity
of $F$ with respect to $p$ and $z$ in (17.53) is dropped).

Recently the Monge-Amp\`ere equation in higher dimensions has flourished
again. A partial differential equations approach to the classical Dirichlet problem
was developed by P.-L. Lions [LP 5, 6] which, on combination with the earlier
considerations of Bakelman, yielded Theorem 17.24 for $C^0(\overline{\Omega}) \cap C^2(\Omega)$ solutions.
Lions' method was based on approximation by problems defined over $\mathbb{R}^n$. In a
similar vein, Cheng and Yau [CY 3, 4] proposed a method based on approxima-
tion by problems with infinite boundary values. Also, a probabilistic approach was
given by Krylov [KV 3]. However the global regularity, which had been an im-
pediment to the direct application of the method of continuity, was finally settled
by Caffarelli, Nirenberg and Spruck [CNS] and Krylov [KV 5, 6], who discovered
Theorems 17.26 and 17.26$'$ respectively, and consequently established Theorem

% PDF page 500
\source{gilbarg-trudinger:page-0500}

\begin{quote}
17.23 in its full strength for the Monge-Amp\`ere equation (17.2). The paper [CNS]
also reduces the solvability of the Dirichlet problem for general Monge-Amp\`ere
equations to the existence of globally smooth subsolutions, from which the case
$\beta = 0$, $\gamma = n$ of Theorem 17.23 follows readily. The work of Caffarelli, Nirenberg
and Spruck [CNS] and also that of Ivochkina [IC 2, 3] embraced the case of
general boundary values $\phi \in C^{4}(\overline{\Omega})$. In [TU], Theorem 17.24 is extended to
boundary data $\partial \Omega \in C^{1,1}$, $\phi \in C^{1,1}(\overline{\Omega})$.

The existence of $C^{2,\alpha}(\overline{\Omega})$ solutions of the Bellman equation, Theorem 17.18$'$, is
due to Krylov [KV 5]. The idea, used in Section 17.8, of invoking the directions
$\gamma$ as variables occurs in Krylov's treatment of interior second derivative estimates
in [KV 4]. Theorem 17.18$'$ may also be extended to cover structure conditions
such as (17.29)$'$, (17.53)$'$ above; (see [TR 13, 14], [CKNS]).

The reduction of nonlinear boundary value problems to a priori estimation in
Theorem 17.30 is due to Fiorenza [FI 2]; (see also Ladyzhenskaya and Ural'tseva
[LU 4]). In our treatment we have adopted some simplification by Lieberman
[LB 2] in the proof of Lemma 17.29, although the main thrust of Lieberman's work
is to replace the method of continuity by another functional analytic approach
which permits a weaker hypothesis than condition (iii) in Theorem 17.28. Oblique
boundary value problems for fully nonlinear equations have been treated recently
by Lions and Trudinger [LPT], [TR 14] for Bellman equations, Lieberman and
Trudinger [LBT], [TR 14] for general nonlinear boundary conditions and by Lions,
Trudinger and Urbas [LTU] for equations of Monge Amp\`ere type.
\end{quote}

\bigskip

\noindent\textbf{Problems}

\medskip

\noindent\textbf{17.1.} \quad Prove that if the operator (17.1) is elliptic at a point $\gamma \in \Gamma$, then the function
$F$ is strictly increasing with respect to $r$ at $\gamma$. That is, there exists a positive number
$\varepsilon$ such that

\[
(17.112)\qquad F(x,z,p,r) < F(x,z,p,r+\eta)
\]

\noindent for all nonzero, positive semidefinite matrices $\eta \in \mathbb{R}^{n\times n}$ with $|\eta|<\varepsilon$.

\medskip

\noindent\textbf{17.2.} \quad Prove that the operator $F$ given by (17.1) is Fr\'echet differentiable as a
mapping from $C^{2,\alpha}(\overline{\Omega})$ into $C^{0,\alpha}(\overline{\Omega})$, for any $\alpha \leq 1$, if the function $F \in C^{2,\alpha}(\overline{\Gamma})$.

\medskip

\noindent\textbf{17.3.} \quad More generally than Problem 17.2, prove that $F$ is Fr\'echet differentiable as
a mapping from $C^{2,\alpha}(\overline{\Omega})$ into $C^{0,\alpha\gamma}(\overline{\Omega})$ for any $\alpha, \beta \leq 1$, $\gamma < \alpha\beta$, if the function $F$ is
differentiable with respect to $z$, $p$ and $r$ with $F$, $F_z$, $F_p$, $F_r \in C^{\beta}(\overline{\Gamma})$.

\medskip

\noindent\textbf{17.4.} \quad Show that Theorem 17.17 remains valid if the condition $F \in C^{2}(\Gamma)$ is
replaced by the existence of the derivatives in (17.53) as weak derivatives in the
sense of Chapter 7.

% PDF page 501
\source{gilbarg-trudinger:page-0501}

\textbf{Problems}

\textbf{17.5.} Let $u \in C^2(\overline{\Omega}) \cap C^4(\Omega)$ be a subharmonic solution of the equation

\begin{equation}
(17.113) \qquad F(D^2u) = g(x)
\end{equation}

in a domain $\Omega \subset \R^n$, where $F \in C^2(\R^{n \times n})$, $g \in C^2(\overline{\Omega})$, $F$ is elliptic with respect to $u$,
tr $[F_{ij}(D^2u)] \ge 1$ and $F$ is concave with respect to $D^2u$. Prove that

\begin{equation}
(17.114) \qquad \sup_{\Omega} |D^2u| \le C\!\left(\sup_{\partial\Omega} |D^2u| + \sup_{\Omega} |D^2g|\right)
\end{equation}

where $C = C(n)$.

\textbf{17.6.} \ (a) Let $F \in C^1(\R^{n \times n})$ be invariant under orthogonal transformations. Show
that for any $r \in \R^{n \times n}$, the matrices $F'(r)$ and $r$ commute.

\hspace{1.7em}(b) Suppose that the conditions of Problem 17.5 are strengthened so that $F$ is
invariant under orthogonal transformations and concave for $r \ge 0$, with det
$[F_{ij}(D^2u)] \ge 1$. Suppose also that the domain $\Omega$ is uniformly convex with boundary
$\partial\Omega \in C^3$ and that $u = 0$ on $\partial\Omega$. Show that if the solution $u$ is convex, then

\begin{equation}
(17.115) \qquad \sup_{\partial\Omega} |D^2u| \le C
\end{equation}

where $C$ depends on $n$, $|g|_{1;\Omega}$, $|u|_{1;\Omega}$ and $\partial\Omega$.

\textbf{17.7.} \ (a) Let $F \in C^1(\R^{n \times n})$ be invariant under orthogonal transformations and
consider for fixed $k,l \in \{1, \dots, n\}$ the polar coordinate transformation

\begin{equation}
(17.116) \qquad x_k = a + r \sin \theta,\qquad x_l = b - r \cos \theta,
\end{equation}

where $a,b$ are constants. If $u \in C^3(\Omega)$ satisfies the equation (17.113) in $\Omega$, show that

\begin{equation}
(17.117) \qquad F_{ij} D_{ij}\!\left(\frac{\partial u}{\partial \theta}\right) = \frac{\partial g}{\partial \theta}
\end{equation}

in $\Omega \cap \{r>0\}$.

\hspace{1.7em}(b) Let $\Omega$ be a $C^3$ domain in $\R^n$. Suppose that $0 \in \partial\Omega$ and that the inner normal
to $\partial\Omega$ at $0$ is directed along the $x_n$ axis. If $u \in C^1(\overline{\Omega})$ and $u = 0$ on $\partial\Omega$, show that in a
neighbourhood of $0$, either

\[
\left|\frac{\partial u}{\partial \theta}\right| \qquad \text{or} \qquad \left|\frac{\partial u}{\partial x_k}\right| \le C|x_k|^2
\]

where $\theta$ is given by (17.116) with appropriate choice of $a,b$ and $l = n$, and $C$ depends
on $\partial\Omega$ and $|u|_{1;\Omega}$. Accordingly show that in Problem 17.6, the assumption of
convexity of $u$ can be dropped. (Cf. [IC 1].)

% PDF page 502
\source{gilbarg-trudinger:page-0502}

17.8.\ Let $\Omega$ be a $C^2$ uniformly convex domain in $\R^n$. Show the existence of a uniformly convex function $u \in C^2(\overline{\Omega})$ vanishing on $\partial \Omega$.

17.9.\ Show that the equation
\[
(17.118)\qquad F[u] = (\Delta u)^2 - |D^2u|^2 = f(x)
\]
is elliptic with respect to any solution $u$ (or its negative) if the function $f$ is positive. Using the results of Problems 17.5, 17.7 together with Theorems 17.14, 17.26, show that the classical Dirichlet problem for (17.118) is solvable for any uniformly convex domain $\Omega \in C^3$, zero boundary values and positive $f \in C^2(\overline{\Omega})$. More generally we note here that this result extends to domains $\Omega$ whose boundary has positive \emph{mean curvature}.

17.10.\ Let $F \in C^1(\Omega \times \R \times \R^n \times \R^{n \times n})$ satisfy
\[
F(x_0,0,0,0)=0,\qquad F_r(x_0,0,0,0)>0
\]
at some point $x_0 \in \Omega$. Using the implicit function (Theorem 17.6), show that in some neighbourhood $\mathcal N$ of $x_0$, there exists a $C^2(\mathcal N)$ solution $u$ of equation (17.1) satisfying $u(x_0)=Du(x_0)=0$. (Hint: Taking $x_0=0$, $t\in\R$, make the change of variables $x=ty$, $u=t^2v$.)


% PDF page 503
\source{gilbarg-trudinger:page-0503}
\begin{thebibliography}{99}

\bibitem{AD}
Adams, R. A.
Sobolev Spaces. New York: Academic Press 1975.

\bibitem{AG}
Agmon, S.
Lectures on Elliptic Boundary Value Problems. Princeton, N. J.: Van Nostrand 1965.

\bibitem{ADN1}
Agmon, S., A. Douglis, and L. Nirenberg
Estimates near the boundary for solutions of elliptic partial differential equations satisfying
general boundary conditions. I. Comm. Pure Appl. Math. \textbf{12}, 623--727 (1959).

\bibitem{ADN2}
Estimates near the boundary for solutions of elliptic partial differential equations satisfying
general boundary conditions. II. Comm. Pure Appl. Math. \textbf{17}, 35--92 (1964).

\bibitem{AL1}
Aleksandrov, A. D.
Dirichlet's problem for the equation Det $\|Z_{ij}\| = \phi$. Vestnik Leningrad Univ. \textbf{13}, no. 1,
5--24 (1958) [Russian].

\bibitem{AL2}
Certain estimates for the Dirichlet problem. Dokl. Akad. Nauk. SSSR \textbf{134}, 1001--1004
(1960) [Russian]. English Translation in Soviet Math. Dokl. \textbf{1}, 1151--1154 (1960).

\bibitem{AL3}
Uniqueness conditions and estimates for the solution of the Dirichlet problem. Vestnik
Leningrad Univ. \textbf{18}, no. 3, 5--29 (1963) [Russian]. English Translation in Amer. Math.
Soc. Transl. (2) \textbf{68}, 89--119 (1968).

\bibitem{AL4}
Majorization of solutions of second-order linear equations. Vestnik Leningrad Univ. \textbf{21},
no. 1, 5--25 (1966) [Russian]. English Translation in Amer. Math. Soc. Transl. (2) \textbf{68},
120--143 (1968).

\bibitem{AL5}
Majorants of solutions and uniqueness conditions for elliptic equations. Vestnik Leningrad
Univ. \textbf{21}, no. 7, 5--20 (1966) [Russian]. English Translation in Amer. Math. Soc. Transl.
(2) \textbf{68}, 144--161 (1968).

\bibitem{AL6}
The impossibility of general estimates for solutions and of uniqueness conditions for linear
equations with norms weaker than in $L_n$. Vestnik Leningrad Univ. \textbf{21}, no. 13, 5--10 (1966)
[Russian]. English Translation in Amer. Math. Soc. Transl. (2) \textbf{68}, 162--168 (1968).

\bibitem{AK}
Alkhutov, Yu. A.
Regularity of boundary points relative to the Dirichlet problem for second order elliptic
equations. Mat. Zametki \textbf{30}, 333--342 (1981) [Russian]. English Translation in Math. Notes
\textbf{30}, 655--661 (1982).

\bibitem{AA}
Allard, W.
On the first variation of a varifold. Ann. of Math. (2) \textbf{95}, 417--491 (1972).

\bibitem{AM}
Almgren, F. J.
Some interior regularity theorems for minimal surfaces and an extension of Bernstein's
theorem. Ann. of Math. (2) \textbf{84}, 277--292 (1966).

\bibitem{AU1}
Aubin, T.
\'{E}quations du type Monge-Amp\`ere sur les vari\'et\'es k\"ahleriennes compactes. C.R. Acad.
Sci. Paris \textbf{283}, 119--121 (1976).

\bibitem{AU2}
Probl\`emes isop\'erim\'etriques et espaces de Sobolev. J. Differential Geometry \textbf{11}, 573--598
(1976).

\bibitem{AU3}
\'{E}quations du type Monge-Amp\`ere sur les vari\'et\'es k\"ahl\'eriennes compactes. Bull. Sci. Math.
\textbf{102}, 63--95 (1978).

\bibitem{AU4}
\'{E}quations de Monge-Amp\`ere r\'eelles. J. Funct. Anal. \textbf{41}, 354--377 (1981).

\bibitem{BA1}
Bakel'man, I. YA.
Generalized solutions of the Monge-Amp\`ere equations. Dokl. Akad. Nauk. SSSR \textbf{114},
1143--1145 (1957) [Russian].

\end{thebibliography}
